% ===================== APPENDIX =====================
% Unlimited pages per REALM CFP. Tables keep their source-order labels tab:tNN.

\section{Multiplicity and inferential scope}\label{app:multiplicity}

\paragraph{Multiplicity status.} The tables in this paper and its appendix are not one inferential family and are not corrected as one. They divide into three classes: confirmatory (hypotheses fixed before the data and corrected within their own family, currently the 2$\times$2 axis-independence set of 64 conjunction hypotheses under BH/Holm \citep{holm1979sequentially} on $\max(p_1, p_2)$); descriptive (intervals, no test, no correction claimed: the full behavioural matrices, the efficiency tables, the failure-attribution counts); and selection-derived (a maximum over layers, signals, thresholds or comparators is part of the statistic, so these bound what a selection could deliver and cannot be read as effect estimates \citep{cawley2010overfitting}; each such table says so in its caption). No family-wise statement covers the set as a whole, and none is claimed.

\section{Deployment classes and ablation}\label{app:classes}

Table~\ref{tab:t02} is the unmatched class comparison (a maximum over four no-image arms, so quote its ordering and never its gaps); Table~\ref{tab:t04} is the class ablation, unmatched and arm-matched; Table~\ref{tab:t05} is the behavioural non-separability tally. It cannot license pooling the four text modes. The tally counts how many metrics separate a mode from the rest in at least seven of eight cells, and under a pure-noise null that bar returns a zero with probability $\approx$76\%, so a mode scoring zero is the expected outcome of no effect and not evidence of sameness. What the table does establish is the positive half: Vision (9 metrics) and SoM (5) are strongly separable from the rest, and that is how the deployment-class grouping of \S\ref{sec:setup} is used: as a convention that survives this evidence rather than one derived from it.

% Table source-order #01: Success rate per mode. Success rate (\%) per observation mode. Denomin
\begin{table}[tp]
\centering
\caption{Success rate per mode. Success rate (\%) per observation mode. Denominator is the canonical scored set (cls 224 / red 203) and, for WebArena, the six-mode task intersection (104). Bold = best mode in the cell. Source: \texttt{representation\_class\_comparison.json}.}
\label{tab:t01}
\footnotesize
\setlength{\tabcolsep}{3pt}
\begin{tabular}{@{}llllllll@{}}
\toprule
cell & DOM & SoM & Vision & DOM+\allowbreak{}stext & DOM+\allowbreak{}sprompt & SoM-image & best \\
\midrule
cls·B0 & 17.41 & \textbf{27.23} & 25.00 & 15.62 & 19.64 & 15.62 & SoM \\
cls·B1 & 6.25 & \textbf{14.29} & 12.50 & 7.59 & 6.70 & 6.70 & SoM \\
cls·B2 & 1.34 & \textbf{2.23} & 2.23 & 0.45 & 1.79 & 0.89 & SoM \\
red·B0 & 14.29 & \textbf{14.78} & 7.39 & 13.30 & 12.32 & 10.84 & SoM \\
red·B1 & 5.91 & \textbf{7.39} & 2.46 & 5.91 & 5.42 & 5.91 & SoM \\
red·B2 & \textbf{3.94} & 0.99 & 1.97 & 1.97 & 0.00 & 0.49 & DOM \\
WA·B0 & 26.92 & 22.12 & 19.23 & \textbf{35.58} & 25.96 & 25.00 & DOM+\allowbreak{}stext \\
WA·B1 & \textbf{16.35} & 13.46 & 9.62 & 16.35 & 16.35 & 11.54 & DOM \\
\bottomrule
\end{tabular}
\end{table}

% Table source-order #02: Best arm per deployment class. Best arm within each deployment class (
\begin{table}[tp]
\centering
\caption{Best arm per deployment class. Best arm within each deployment class (\%). Classes: no-image = \{DOM, DOM+stext, DOM+sprompt, SoM-image\}, vision-only = \{Vision\}, hybrid = \{SoM\}. Grouping the four is a deployment convention, not an equivalence claim: over 26 metrics none of them is the repeated extreme (Table~\ref{tab:t05}), which is a counting diagnostic rather than a separability test. Source: \texttt{representation\_class\_comparison.json}.}
\label{tab:t02}
\footnotesize
\setlength{\tabcolsep}{3pt}
\begin{tabular}{@{}lllll@{}}
\toprule
cell & no-image (4 arms) & vision-only & hybrid & sole best \\
\midrule
cls·B0 & 19.64 (DOM+\allowbreak{}sprompt) & 25.00 & 27.23 & hybrid \\
cls·B1 & 7.59 (DOM+\allowbreak{}stext) & 12.50 & 14.29 & hybrid \\
cls·B2 & 1.79 (DOM+\allowbreak{}sprompt) & 2.23 & 2.23 & tie: hybrid+vision-only \\
red·B0 & 14.29 (DOM) & 7.39 & 14.78 & hybrid \\
red·B1 & 5.91 (DOM) & 2.46 & 7.39 & hybrid \\
red·B2 & 3.94 (DOM) & 1.97 & 0.99 & no-image \\
WA·B0 & 35.58 (DOM+\allowbreak{}stext) & 19.23 & 22.12 & no-image \\
WA·B1 & 16.35 (DOM) & 9.62 & 13.46 & no-image \\
\bottomrule
\end{tabular}
\end{table}

% Table source-order #03: Deployment classes at one arm each. The same comparison at \textbf{one
\begin{table}[tp]
\centering
\caption{Deployment classes at one arm each. The same comparison at one arm per class (\%). Table~\ref{tab:t02}'s no-image column is a maximum over four arms while the other two are single arms, so it is biased up; this panel uses the arm of each class that exists outside this study. The ordering does not move: hybrid 4, no-image 3, one tie, and vision-only is never a sole best in either version: which is the robustness statement Table~\ref{tab:t02} cannot make. The gaps do move: on \texttt{WA·B0} the no-image lead is +4.81pp here against +13.46pp there, because Table~\ref{tab:t02}'s figure is carried by DOM+stext. Quote this table for any class gap; Table~\ref{tab:t02} only for the ordering. Source: \texttt{representation\_class\_comparison.json}.}
\label{tab:t03}
\footnotesize
\setlength{\tabcolsep}{3pt}
\begin{tabular}{@{}
  >{\raggedright\arraybackslash}p{(\linewidth - 12\tabcolsep) * \real{0.1429}}
  >{\raggedright\arraybackslash}p{(\linewidth - 12\tabcolsep) * \real{0.1429}}
  >{\raggedright\arraybackslash}p{(\linewidth - 12\tabcolsep) * \real{0.1429}}
  >{\raggedright\arraybackslash}p{(\linewidth - 12\tabcolsep) * \real{0.1429}}
  >{\raggedright\arraybackslash}p{(\linewidth - 12\tabcolsep) * \real{0.1429}}
  >{\raggedright\arraybackslash}p{(\linewidth - 12\tabcolsep) * \real{0.1429}}
  >{\raggedright\arraybackslash}p{(\linewidth - 12\tabcolsep) * \real{0.1429}}@{}}
\toprule
\begin{minipage}[b]{\linewidth}\raggedright
cell
\end{minipage} & \begin{minipage}[b]{\linewidth}\raggedright
no-image (DOM)
\end{minipage} & \begin{minipage}[b]{\linewidth}\raggedright
vision-only (Vision)
\end{minipage} & \begin{minipage}[b]{\linewidth}\raggedright
hybrid (SoM)
\end{minipage} & \begin{minipage}[b]{\linewidth}\raggedright
sole best
\end{minipage} & \begin{minipage}[b]{\linewidth}\raggedright
gap vs hybrid
\end{minipage} & \begin{minipage}[b]{\linewidth}\raggedright
Table 2 gap
\end{minipage} \\
\midrule
cls·B0 & 17.41 & 25.00 & 27.23 & hybrid & -9.82 & -7.59 \\
cls·B1 & 6.25 & 12.50 & 14.29 & hybrid & -8.04 & -6.70 \\
cls·B2 & 1.34 & 2.23 & 2.23 & tie: hybrid+vision-only & -0.89 & -0.45 \\
red·B0 & 14.29 & 7.39 & 14.78 & hybrid & -0.49 & -0.49 \\
red·B1 & 5.91 & 2.46 & 7.39 & hybrid & -1.48 & -1.48 \\
red·B2 & 3.94 & 1.97 & 0.99 & no-image & +2.96 & +2.96 \\
WA·B0 & 26.92 & 19.23 & 22.12 & no-image & +4.81 & +13.46 \\
WA·B1 & 16.35 & 9.62 & 13.46 & no-image & +2.88 & +2.88 \\
\bottomrule
\end{tabular}
\end{table}

% Table source-order #04: Class ablation, unmatched and arm-matched. Class ablation. Columns 2--
\begin{table}[tp]
\centering
\caption{Class ablation, unmatched and arm-matched. Class ablation. Columns 2--4: oracle coverage lost when a whole class is unavailable: not arm-matched, the no-image class has four arms and the others one each, so most of the gap is arm count. Columns 5--7: the matched comparison: gain from adding ONE arm of that class to the cell's best single arm (: = that class already supplies the starting arm). The matched panel shows no systematic difference between classes. Source: \texttt{representation\_class\_comparison.json}.}
\label{tab:t04}
\footnotesize
\setlength{\tabcolsep}{3pt}
\begin{tabular}{@{}
  >{\raggedright\arraybackslash}p{(\linewidth - 14\tabcolsep) * \real{0.1250}}
  >{\raggedright\arraybackslash}p{(\linewidth - 14\tabcolsep) * \real{0.1250}}
  >{\raggedright\arraybackslash}p{(\linewidth - 14\tabcolsep) * \real{0.1250}}
  >{\raggedright\arraybackslash}p{(\linewidth - 14\tabcolsep) * \real{0.1250}}
  >{\raggedright\arraybackslash}p{(\linewidth - 14\tabcolsep) * \real{0.1250}}
  >{\raggedright\arraybackslash}p{(\linewidth - 14\tabcolsep) * \real{0.1250}}
  >{\raggedright\arraybackslash}p{(\linewidth - 14\tabcolsep) * \real{0.1250}}
  >{\raggedright\arraybackslash}p{(\linewidth - 14\tabcolsep) * \real{0.1250}}@{}}
\toprule
\begin{minipage}[b]{\linewidth}\raggedright
cell
\end{minipage} & \begin{minipage}[b]{\linewidth}\raggedright
all six
\end{minipage} & \begin{minipage}[b]{\linewidth}\raggedright
−no-image
\end{minipage} & \begin{minipage}[b]{\linewidth}\raggedright
−vision-only
\end{minipage} & \begin{minipage}[b]{\linewidth}\raggedright
−hybrid
\end{minipage} & \begin{minipage}[b]{\linewidth}\raggedright
+1 no-image
\end{minipage} & \begin{minipage}[b]{\linewidth}\raggedright
+1 vision-only
\end{minipage} & \begin{minipage}[b]{\linewidth}\raggedright
+1 hybrid
\end{minipage} \\
\midrule
cls·B0 & 43.30 & −9.38 & −4.02 & −2.68 & +7.14pp & +6.70pp & --- \\
cls·B1 & 24.55 & −5.36 & −4.02 & −4.46 & +3.57pp & +4.91pp & --- \\
cls·B2 & 7.14 & −2.68 & −2.23 & −2.23 & +1.79pp & +2.23pp & --- \\
red·B0 & 26.11 & −7.88 & −1.97 & −1.97 & +4.93pp & +3.45pp & --- \\
red·B1 & 11.82 & −3.45 & −0.99 & −0.49 & +1.97pp & +0.99pp & --- \\
red·B2 & 7.39 & −4.43 & −1.48 & −0.49 & +0.99pp & +1.97pp & +0.49pp \\
WA·B0 & 51.92 & −22.12 & −3.85 & −1.92 & +5.77pp & +5.77pp & +4.81pp \\
WA·B1 & 30.77 & −14.42 & −0.96 & −0.96 & +4.81pp & +3.85pp & +3.85pp \\
\bottomrule
\end{tabular}
\end{table}

% Table source-order #05: Absence of repeated extrema among image-free modes. \textbf{Absence of
\begin{table}[tp]
\centering
\caption{Absence of repeated extrema among image-free modes. Absence of repeated extrema: read the name literally. A mode `reaches the bar' on a metric when it is the extreme (highest or lowest) in ≥7 of 8 cells (87.5\%). Over 26 metrics the four image-free modes reach it on none. \textbf{Caution:} This is not a separability test. \texttt{rank\_consistency()} only counts which mode attains each metric's max/min; a mode that is consistently \emph{second} (and strongly distinguishable) never reaches the bar. Establishing non-separability would need pairwise equivalence margins with task-clustered intervals, which this does not do. \textbf{Caution:} The threshold got stricter when cells grew, and an earlier caption said otherwise: ≥7/8 is 87.5\%, not the 83\% it claimed, and the six-cell ≥5/6 it says it matches is 83.3\%: a +4.2pp shift, chosen after the cell count changed. Carrying the literal numerator (≥5/8 = 62.5\%) instead would let DOM+stext clear two metrics and this negative would appear to break, so the choice is load-bearing and is disclosed rather than defended. Source: \texttt{per\_mode\_four\_dimension\_profile\_with\_wa.json}}
\label{tab:t05}
\footnotesize
\setlength{\tabcolsep}{3pt}
\begin{tabular}{@{}ll@{}}
\toprule
mode & metrics reaching the bar \\
\midrule
Vision & 9 \\
SoM & 5 \\
DOM & 0 \\
DOM+\allowbreak{}stext & 0 \\
DOM+\allowbreak{}sprompt & 0 \\
SoM-image & 0 \\
\bottomrule
\end{tabular}
\end{table}


\section{Full behavioural matrices}\label{app:matrices}

The four dimensions unsummarised, for every metric in every cell and mode: Outcome, Macro (what the agent did), Micro (how often what it did failed), and Efficiency. These are the substrate the consistency counts in Table~\ref{tab:t05} are computed from, included so a reader can check a claim rather than take the tally on trust. Two columns are gates rather than measurements, and the tables mark them: \texttt{loc-fallback} is near-zero for Vision because there are no element ids to fall back from, not because it fails less.

% Table source-order #06: Full matrix --- Outcome dimension. Outcome dimension, every cell × eve
\begin{table}[tp]
\centering
\caption{Full matrix: Outcome dimension. Outcome dimension, every cell × every mode. \texttt{unique} counts tasks no other mode in that cell solved. Denominators: cls 224 / red 203 / WA 104. Source: \texttt{per\_mode\_four\_dimension\_profile\_with\_wa.json}.}
\label{tab:t06}
\footnotesize
\setlength{\tabcolsep}{3pt}
\begin{tabular}{@{}lllll@{}}
\toprule
cell & mode & SR \% & solves & unique \\
\midrule
cls·B0 & DOM & 17.41 & 39 & 4 \\
cls·B0 & SoM & 27.23 & 61 & 6 \\
cls·B0 & Vision & 25.00 & 56 & 9 \\
cls·B0 & DOM+\allowbreak{}stext & 15.62 & 35 & 2 \\
cls·B0 & DOM+\allowbreak{}sprompt & 19.64 & 44 & 6 \\
cls·B0 & SoM-image & 15.62 & 35 & 2 \\
cls·B1 & DOM & 6.25 & 14 & 1 \\
cls·B1 & SoM & 14.29 & 32 & 10 \\
cls·B1 & Vision & 12.50 & 28 & 9 \\
cls·B1 & DOM+\allowbreak{}stext & 7.59 & 17 & 1 \\
cls·B1 & DOM+\allowbreak{}sprompt & 6.70 & 15 & 2 \\
cls·B1 & SoM-image & 6.70 & 15 & 3 \\
cls·B2 & DOM & 1.34 & 3 & 1 \\
cls·B2 & SoM & 2.23 & 5 & 5 \\
cls·B2 & Vision & 2.23 & 5 & 5 \\
cls·B2 & DOM+\allowbreak{}stext & 0.45 & 1 & 0 \\
cls·B2 & DOM+\allowbreak{}sprompt & 1.79 & 4 & 0 \\
cls·B2 & SoM-image & 0.89 & 2 & 1 \\
red·B0 & DOM & 14.29 & 29 & 3 \\
red·B0 & SoM & 14.78 & 30 & 4 \\
red·B0 & Vision & 7.39 & 15 & 4 \\
red·B0 & DOM+\allowbreak{}stext & 13.30 & 27 & 2 \\
red·B0 & DOM+\allowbreak{}sprompt & 12.32 & 25 & 2 \\
red·B0 & SoM-image & 10.84 & 22 & 2 \\
red·B1 & DOM & 5.91 & 12 & 1 \\
red·B1 & SoM & 7.39 & 15 & 1 \\
red·B1 & Vision & 2.46 & 5 & 2 \\
red·B1 & DOM+\allowbreak{}stext & 5.91 & 12 & 1 \\
red·B1 & DOM+\allowbreak{}sprompt & 5.42 & 11 & 2 \\
red·B1 & SoM-image & 5.91 & 12 & 0 \\
red·B2 & DOM & 3.94 & 8 & 6 \\
red·B2 & SoM & 0.99 & 2 & 1 \\
red·B2 & Vision & 1.97 & 4 & 3 \\
red·B2 & DOM+\allowbreak{}stext & 1.97 & 4 & 1 \\
red·B2 & DOM+\allowbreak{}sprompt & 0.00 & 0 & 0 \\
red·B2 & SoM-image & 0.49 & 1 & 1 \\
WA·B0 & DOM & 26.92 & 28 & 2 \\
WA·B0 & SoM & 22.12 & 23 & 2 \\
WA·B0 & Vision & 19.23 & 20 & 4 \\
WA·B0 & DOM+\allowbreak{}stext & 35.58 & 37 & 7 \\
WA·B0 & DOM+\allowbreak{}sprompt & 25.96 & 27 & 2 \\
WA·B0 & SoM-image & 25.00 & 26 & 1 \\
WA·B1 & DOM & 16.35 & 17 & 3 \\
WA·B1 & SoM & 13.46 & 14 & 1 \\
WA·B1 & Vision & 9.62 & 10 & 1 \\
WA·B1 & DOM+\allowbreak{}stext & 16.35 & 17 & 3 \\
WA·B1 & DOM+\allowbreak{}sprompt & 16.35 & 17 & 5 \\
WA·B1 & SoM-image & 11.54 & 12 & 1 \\
\bottomrule
\end{tabular}
\end{table}

% Table source-order #07: Full matrix --- Macro dimension. Macro dimension --- what the agent di
\begin{table}[tp]
\centering
\caption{Full matrix (Macro dimension. Macro dimension) what the agent did, per step, aggregated per episode. Fractions are over agent actions. \texttt{cap-hit} = share of episodes that exhausted the 30-step budget. Source: \texttt{per\_mode\_four\_dimension\_profile\_with\_wa.json}.}
\label{tab:t07}
\scriptsize
\setlength{\tabcolsep}{3pt}
\begin{tabular}{@{}lllllllll@{}}
\toprule
cell & mode & steps/ep & cap-hit & click & type & scroll & search-loop & URL-revisit \\
\midrule
cls·B0 & DOM & 15.61 & 0.268 & 0.322 & 0.224 & 0.183 & 0.812 & 0.606 \\
cls·B0 & SoM & 13.67 & 0.259 & 0.319 & 0.205 & 0.155 & 0.692 & 0.558 \\
cls·B0 & Vision & 15.88 & 0.295 & 0.353 & 0.146 & 0.260 & 0.728 & 0.639 \\
cls·B0 & DOM+stext & 15.83 & 0.295 & 0.322 & 0.194 & 0.196 & 0.808 & 0.603 \\
cls·B0 & DOM+sprompt & 14.96 & 0.290 & 0.328 & 0.208 & 0.174 & 0.759 & 0.583 \\
cls·B0 & SoM-image & 16.23 & 0.321 & 0.312 & 0.201 & 0.208 & 0.790 & 0.600 \\
cls·B1 & DOM & 21.38 & 0.585 & 0.246 & 0.345 & 0.171 & 0.777 & 0.698 \\
cls·B1 & SoM & 18.01 & 0.491 & 0.368 & 0.211 & 0.137 & 0.643 & 0.655 \\
cls·B1 & Vision & 20.17 & 0.598 & 0.330 & 0.071 & 0.360 & 0.603 & 0.745 \\
cls·B1 & DOM+stext & 22.46 & 0.625 & 0.281 & 0.367 & 0.148 & 0.804 & 0.716 \\
cls·B1 & DOM+sprompt & 21.40 & 0.594 & 0.375 & 0.209 & 0.163 & 0.723 & 0.710 \\
cls·B1 & SoM-image & 21.26 & 0.576 & 0.436 & 0.212 & 0.154 & 0.737 & 0.715 \\
cls·B2 & DOM & 27.38 & 0.817 & 0.520 & 0.142 & 0.031 & 0.647 & 0.823 \\
cls·B2 & SoM & 24.37 & 0.688 & 0.488 & 0.143 & 0.033 & 0.616 & 0.828 \\
cls·B2 & Vision & 28.25 & 0.888 & 0.336 & 0.128 & 0.323 & 0.710 & 0.896 \\
cls·B2 & DOM+stext & 26.85 & 0.804 & 0.384 & 0.168 & 0.043 & 0.629 & 0.832 \\
cls·B2 & DOM+sprompt & 27.84 & 0.830 & 0.514 & 0.078 & 0.048 & 0.629 & 0.857 \\
cls·B2 & SoM-image & 28.38 & 0.853 & 0.523 & 0.082 & 0.033 & 0.683 & 0.858 \\
red·B0 & DOM & 20.18 & 0.498 & 0.455 & 0.148 & 0.166 & 0.473 & 0.721 \\
red·B0 & SoM & 20.08 & 0.512 & 0.474 & 0.123 & 0.099 & 0.363 & 0.731 \\
red·B0 & Vision & 23.55 & 0.672 & 0.339 & 0.079 & 0.343 & 0.239 & 0.817 \\
red·B0 & DOM+stext & 23.22 & 0.647 & 0.442 & 0.153 & 0.117 & 0.388 & 0.775 \\
red·B0 & DOM+sprompt & 19.87 & 0.463 & 0.475 & 0.142 & 0.126 & 0.428 & 0.704 \\
red·B0 & SoM-image & 22.90 & 0.592 & 0.453 & 0.137 & 0.124 & 0.338 & 0.775 \\
red·B1 & DOM & 23.44 & 0.680 & 0.471 & 0.238 & 0.066 & 0.596 & 0.754 \\
red·B1 & SoM & 22.38 & 0.670 & 0.574 & 0.129 & 0.095 & 0.409 & 0.747 \\
red·B1 & Vision & 23.24 & 0.685 & 0.434 & 0.029 & 0.270 & 0.148 & 0.812 \\
red·B1 & DOM+stext & 25.56 & 0.783 & 0.476 & 0.246 & 0.067 & 0.606 & 0.793 \\
red·B1 & DOM+sprompt & 23.72 & 0.695 & 0.545 & 0.149 & 0.052 & 0.567 & 0.757 \\
red·B1 & SoM-image & 25.17 & 0.773 & 0.580 & 0.166 & 0.056 & 0.557 & 0.778 \\
red·B2 & DOM & 28.39 & 0.897 & 0.724 & 0.108 & 0.042 & 0.227 & 0.870 \\
red·B2 & SoM & 26.34 & 0.783 & 0.702 & 0.103 & 0.043 & 0.148 & 0.857 \\
red·B2 & Vision & 26.91 & 0.828 & 0.349 & 0.081 & 0.344 & 0.138 & 0.911 \\
red·B2 & DOM+stext & 27.27 & 0.818 & 0.660 & 0.099 & 0.038 & 0.266 & 0.885 \\
red·B2 & DOM+sprompt & 27.68 & 0.842 & 0.703 & 0.081 & 0.052 & 0.271 & 0.873 \\
red·B2 & SoM-image & 27.87 & 0.882 & 0.708 & 0.094 & 0.048 & 0.217 & 0.882 \\
WA·B0 & DOM & 16.88 & 0.298 & 0.544 & 0.208 & 0.072 & 0.308 & 0.709 \\
WA·B0 & SoM & 17.45 & 0.365 & 0.465 & 0.278 & 0.029 & 0.327 & 0.716 \\
WA·B0 & Vision & 22.38 & 0.567 & 0.409 & 0.200 & 0.241 & 0.240 & 0.783 \\
WA·B0 & DOM+stext & 19.77 & 0.471 & 0.478 & 0.323 & 0.045 & 0.288 & 0.741 \\
WA·B0 & DOM+sprompt & 17.08 & 0.317 & 0.538 & 0.198 & 0.068 & 0.260 & 0.714 \\
WA·B0 & SoM-image & 18.97 & 0.385 & 0.502 & 0.284 & 0.056 & 0.221 & 0.729 \\
WA·B1 & DOM & 22.64 & 0.625 & 0.459 & 0.319 & 0.023 & 0.644 & 0.746 \\
WA·B1 & SoM & 23.91 & 0.702 & 0.499 & 0.218 & 0.035 & 0.452 & 0.793 \\
WA·B1 & Vision & 23.12 & 0.606 & 0.450 & 0.097 & 0.249 & 0.260 & 0.807 \\
WA·B1 & DOM+stext & 23.33 & 0.615 & 0.495 & 0.316 & 0.028 & 0.683 & 0.733 \\
WA·B1 & DOM+sprompt & 24.50 & 0.702 & 0.551 & 0.231 & 0.033 & 0.519 & 0.776 \\
WA·B1 & SoM-image & 23.79 & 0.683 & 0.569 & 0.249 & 0.017 & 0.587 & 0.772 \\
\bottomrule
\end{tabular}
\end{table}

% Table source-order #08: Full matrix --- Micro dimension. Micro dimension --- per-step executio
\begin{table}[tp]
\centering
\caption{Full matrix (Micro dimension. Micro dimension) per-step execution quality. \texttt{act-fail\textbar{}click} and \texttt{act-fail\textbar{}type} are conditional on the action type, so they are not comparable to the unconditional \texttt{act-fail} column. \textbf{Caution:} They are also not pure conditionals. An episode containing no click at all has no click-failure rate to report, and the producer stores \texttt{0.0} there (an undefined rate encoded as a perfect one) which is then averaged over every task. The zero-denominator share is large: 25--35\% of episodes never type, and a few percent never click. The product now carries \texttt{*\_fail\_rate\_complete\_case} (averaged only over episodes where the action occurred) and \texttt{*\_fail\_rate\_denom\_zero\_frac} beside these columns; the complete-case values run higher, and any statement about which mode fails most on clicks or types should be read from those, not from this column. \textbf{Caution:} \texttt{loc-fallback} is near-zero for Vision by construction (no element ids to fall back from), not as a finding. Source: \texttt{per\_mode\_four\_dimension\_profile\_with\_wa.json}.}
\label{tab:t08}
\scriptsize
\setlength{\tabcolsep}{3pt}
\begin{tabular}{@{}lllllllllllll@{}}
\toprule
cell & mode & parse-fail & act-fail & act-fail & click & act-fail & type & no-op & scroll-inert & no-op & success & vis-gap \\
\midrule
cls·B0 & DOM & 0.0007 & 0.133 & 0.163 & 0.035 & 0.241 & 0.096 & 0.108 & 0.046 & 0.078 & 0.324 & 0.737 \\
cls·B0 & SoM & 0.0021 & 0.082 & 0.102 & 0.007 & 0.216 & 0.018 & 0.134 & 0.051 & 0.027 & 0.300 & 0.741 \\
cls·B0 & Vision & 0.0000 & 0.150 & 0.153 & 0.005 & 0.263 & 0.072 & 0.113 & 0.058 & 0.002 & 0.396 & 0.705 \\
cls·B0 & DOM+stext & 0.0000 & 0.101 & 0.100 & 0.022 & 0.209 & 0.071 & 0.108 & 0.042 & 0.046 & 0.310 & 0.705 \\
cls·B0 & DOM+sprompt & 0.0009 & 0.133 & 0.159 & 0.026 & 0.246 & 0.078 & 0.114 & 0.049 & 0.079 & 0.322 & 0.719 \\
cls·B0 & SoM-image & 0.0005 & 0.112 & 0.130 & 0.016 & 0.221 & 0.078 & 0.109 & 0.042 & 0.037 & 0.332 & 0.688 \\
cls·B1 & DOM & 0.0031 & 0.285 & 0.155 & 0.153 & 0.359 & 0.219 & 0.074 & 0.011 & 0.175 & 0.443 & 0.411 \\
cls·B1 & SoM & 0.0016 & 0.249 & 0.203 & 0.068 & 0.353 & 0.149 & 0.104 & 0.013 & 0.118 & 0.390 & 0.509 \\
cls·B1 & Vision & 0.0000 & 0.454 & 0.345 & 0.031 & 0.548 & 0.278 & 0.094 & 0.018 & 0.011 & 0.530 & 0.402 \\
cls·B1 & DOM+stext & 0.0004 & 0.259 & 0.182 & 0.140 & 0.317 & 0.159 & 0.058 & 0.010 & 0.202 & 0.461 & 0.384 \\
cls·B1 & DOM+sprompt & 0.0044 & 0.309 & 0.272 & 0.081 & 0.380 & 0.178 & 0.070 & 0.012 & 0.174 & 0.409 & 0.406 \\
cls·B1 & SoM-image & 0.0010 & 0.322 & 0.274 & 0.136 & 0.381 & 0.169 & 0.059 & 0.015 & 0.203 & 0.461 & 0.424 \\
cls·B2 & DOM & 0.0259 & 0.509 & 0.555 & 0.207 & 0.526 & 0.099 & 0.018 & 0.030 & 0.339 & 0.630 & 0.125 \\
cls·B2 & SoM & 0.0876 & 0.559 & 0.475 & 0.110 & 0.584 & 0.067 & 0.025 & 0.026 & 0.327 & 0.627 & 0.152 \\
cls·B2 & Vision & 0.0067 & 0.670 & 0.773 & 0.355 & 0.685 & 0.443 & 0.015 & 0.017 & 0.124 & 0.518 & 0.103 \\
cls·B2 & DOM+stext & 0.0761 & 0.441 & 0.438 & 0.227 & 0.453 & 0.130 & 0.012 & 0.029 & 0.292 & 0.602 & 0.062 \\
cls·B2 & DOM+sprompt & 0.0518 & 0.495 & 0.569 & 0.167 & 0.504 & 0.137 & 0.009 & 0.028 & 0.312 & 0.633 & 0.071 \\
cls·B2 & SoM-image & 0.0522 & 0.425 & 0.507 & 0.109 & 0.433 & 0.099 & 0.008 & 0.027 & 0.273 & 0.649 & 0.054 \\
red·B0 & DOM & 0.0002 & 0.222 & 0.194 & 0.051 & 0.294 & 0.161 & 0.072 & 0.075 & 0.119 & 0.428 & 0.507 \\
red·B0 & SoM & 0.0008 & 0.252 & 0.127 & 0.022 & 0.328 & 0.085 & 0.077 & 0.092 & 0.068 & 0.491 & 0.488 \\
red·B0 & Vision & 0.0000 & 0.382 & 0.193 & 0.018 & 0.430 & 0.268 & 0.049 & 0.064 & 0.006 & 0.556 & 0.328 \\
red·B0 & DOM+stext & 0.0005 & 0.292 & 0.126 & 0.062 & 0.339 & 0.146 & 0.047 & 0.087 & 0.064 & 0.518 & 0.358 \\
red·B0 & DOM+sprompt & 0.0026 & 0.194 & 0.146 & 0.046 & 0.270 & 0.149 & 0.076 & 0.093 & 0.083 & 0.421 & 0.537 \\
red·B0 & SoM-image & 0.0027 & 0.285 & 0.118 & 0.034 & 0.335 & 0.160 & 0.050 & 0.097 & 0.070 & 0.502 & 0.408 \\
red·B1 & DOM & 0.0021 & 0.227 & 0.210 & 0.049 & 0.274 & 0.090 & 0.047 & 0.050 & 0.143 & 0.465 & 0.315 \\
red·B1 & SoM & 0.0063 & 0.297 & 0.228 & 0.062 & 0.357 & 0.111 & 0.061 & 0.061 & 0.176 & 0.552 & 0.320 \\
red·B1 & Vision & 0.0141 & 0.532 & 0.334 & 0.010 & 0.582 & 0.302 & 0.050 & 0.022 & 0.004 & 0.632 & 0.291 \\
red·B1 & DOM+stext & 0.0018 & 0.283 & 0.195 & 0.107 & 0.319 & 0.156 & 0.036 & 0.042 & 0.164 & 0.511 & 0.217 \\
red·B1 & DOM+sprompt & 0.0010 & 0.295 & 0.278 & 0.056 & 0.340 & 0.114 & 0.045 & 0.055 & 0.201 & 0.494 & 0.305 \\
red·B1 & SoM-image & 0.0019 & 0.332 & 0.297 & 0.077 & 0.368 & 0.146 & 0.036 & 0.059 & 0.208 & 0.544 & 0.222 \\
red·B2 & DOM & 0.0224 & 0.492 & 0.502 & 0.111 & 0.501 & 0.156 & 0.009 & 0.044 & 0.443 & 0.770 & 0.074 \\
red·B2 & SoM & 0.0431 & 0.384 & 0.344 & 0.090 & 0.403 & 0.053 & 0.019 & 0.052 & 0.305 & 0.760 & 0.163 \\
red·B2 & Vision & 0.0163 & 0.640 & 0.613 & 0.276 & 0.669 & 0.547 & 0.029 & 0.005 & 0.078 & 0.597 & 0.153 \\
red·B2 & DOM+stext & 0.0786 & 0.285 & 0.263 & 0.081 & 0.292 & 0.094 & 0.007 & 0.041 & 0.202 & 0.728 & 0.059 \\
red·B2 & DOM+sprompt & 0.0660 & 0.604 & 0.621 & 0.126 & 0.609 & 0.115 & 0.005 & 0.042 & 0.503 & 0.724 & 0.044 \\
red·B2 & SoM-image & 0.0501 & 0.316 & 0.297 & 0.123 & 0.325 & 0.137 & 0.009 & 0.034 & 0.240 & 0.725 & 0.049 \\
WA·B0 & DOM & 0.0008 & 0.277 & 0.286 & 0.083 & 0.363 & 0.046 & 0.085 & 0.226 & 0.165 & 0.455 & 0.702 \\
WA·B0 & SoM & 0.0063 & 0.249 & 0.139 & 0.079 & 0.345 & 0.026 & 0.096 & 0.219 & 0.045 & 0.433 & 0.625 \\
WA·B0 & Vision & 0.0000 & 0.245 & 0.178 & 0.042 & 0.295 & 0.165 & 0.050 & 0.152 & 0.007 & 0.478 & 0.433 \\
WA·B0 & DOM+stext & 0.0010 & 0.347 & 0.292 & 0.200 & 0.421 & 0.057 & 0.074 & 0.247 & 0.099 & 0.497 & 0.538 \\
WA·B0 & DOM+sprompt & 0.0018 & 0.295 & 0.278 & 0.070 & 0.379 & 0.045 & 0.085 & 0.223 & 0.157 & 0.469 & 0.692 \\
WA·B0 & SoM-image & 0.0033 & 0.302 & 0.268 & 0.170 & 0.377 & 0.040 & 0.075 & 0.268 & 0.092 & 0.500 & 0.625 \\
WA·B1 & DOM & 0.0227 & 0.244 & 0.201 & 0.038 & 0.289 & 0.026 & 0.045 & 0.135 & 0.139 & 0.512 & 0.327 \\
WA·B1 & SoM & 0.0052 & 0.363 & 0.214 & 0.084 & 0.415 & 0.024 & 0.052 & 0.166 & 0.149 & 0.573 & 0.279 \\
WA·B1 & Vision & 0.0072 & 0.416 & 0.326 & 0.023 & 0.456 & 0.215 & 0.041 & 0.102 & 0.002 & 0.571 & 0.365 \\
WA·B1 & DOM+stext & 0.0102 & 0.206 & 0.197 & 0.076 & 0.247 & 0.045 & 0.041 & 0.113 & 0.105 & 0.503 & 0.385 \\
WA·B1 & DOM+sprompt & 0.0068 & 0.343 & 0.302 & 0.070 & 0.379 & 0.066 & 0.036 & 0.098 & 0.197 & 0.522 & 0.298 \\
WA·B1 & SoM-image & 0.0144 & 0.343 & 0.315 & 0.076 & 0.382 & 0.036 & 0.039 & 0.099 & 0.180 & 0.562 & 0.288 \\
\bottomrule
\end{tabular}
\end{table}

% Table source-order #09: Full matrix --- Efficiency dimension. Efficiency dimension. \textbf{Co
\begin{table}[tp]
\centering
\caption{Full matrix (Efficiency dimension. Efficiency dimension. Cost is comparable within a cell only) B0 bills a proxy API, B1/B2 are electricity-derived from a per-token constant calibrated for a different accelerator, so absolute dollars for B1/B2 are uncalibrated and only \texttt{cost\ rel\ DOM} is safe across backbones. \texttt{latency\ canon} removes retry, busy-wait and recovered-screenshot time; it differs from raw only on the API-served arm. Source: \texttt{per\_mode\_four\_dimension\_profile\_with\_wa.json}.}
\label{tab:t09}
\footnotesize
\setlength{\tabcolsep}{3pt}
\begin{tabular}{@{}
  >{\raggedright\arraybackslash}p{(\linewidth - 12\tabcolsep) * \real{0.1429}}
  >{\raggedright\arraybackslash}p{(\linewidth - 12\tabcolsep) * \real{0.1429}}
  >{\raggedright\arraybackslash}p{(\linewidth - 12\tabcolsep) * \real{0.1429}}
  >{\raggedright\arraybackslash}p{(\linewidth - 12\tabcolsep) * \real{0.1429}}
  >{\raggedright\arraybackslash}p{(\linewidth - 12\tabcolsep) * \real{0.1429}}
  >{\raggedright\arraybackslash}p{(\linewidth - 12\tabcolsep) * \real{0.1429}}
  >{\raggedright\arraybackslash}p{(\linewidth - 12\tabcolsep) * \real{0.1429}}@{}}
\toprule
\begin{minipage}[b]{\linewidth}\raggedright
cell
\end{minipage} & \begin{minipage}[b]{\linewidth}\raggedright
mode
\end{minipage} & \begin{minipage}[b]{\linewidth}\raggedright
cost/ep
\end{minipage} & \begin{minipage}[b]{\linewidth}\raggedright
cost rel DOM
\end{minipage} & \begin{minipage}[b]{\linewidth}\raggedright
latency s
\end{minipage} & \begin{minipage}[b]{\linewidth}\raggedright
latency canon s
\end{minipage} & \begin{minipage}[b]{\linewidth}\raggedright
tokens/ep
\end{minipage} \\
\midrule
cls·B0 & DOM & 0.06962 & 1.000 & 115.0 & 114.1 & 63045 \\
cls·B0 & SoM & 0.07236 & 1.039 & 106.7 & 106.0 & 67000 \\
cls·B0 & Vision & 0.06481 & 0.931 & 126.3 & 125.2 & 58302 \\
cls·B0 & DOM+\allowbreak{}stext & 0.06919 & 0.994 & 123.8 & 120.4 & 62154 \\
cls·B0 & DOM+\allowbreak{}sprompt & 0.06853 & 0.984 & 109.4 & 107.8 & 62555 \\
cls·B0 & SoM-image & 0.07206 & 1.035 & 121.1 & 117.9 & 65334 \\
cls·B1 & DOM & 0.05951 & 1.000 & 308.9 & 308.9 & 61904 \\
cls·B1 & SoM & 0.06028 & 1.013 & 262.0 & 262.0 & 62953 \\
cls·B1 & Vision & 0.04316 & 0.725 & 269.8 & 269.8 & 44524 \\
cls·B1 & DOM+\allowbreak{}stext & 0.05879 & 0.988 & 313.5 & 313.5 & 61074 \\
cls·B1 & DOM+\allowbreak{}sprompt & 0.06304 & 1.059 & 301.4 & 301.4 & 65628 \\
cls·B1 & SoM-image & 0.05970 & 1.003 & 311.7 & 311.7 & 61985 \\
cls·B2 & DOM & 0.07676 & 1.000 & 402.3 & 402.3 & 79653 \\
cls·B2 & SoM & 0.09075 & 1.182 & 374.3 & 374.3 & 95081 \\
cls·B2 & Vision & 0.07065 & 0.920 & 417.8 & 417.8 & 73126 \\
cls·B2 & DOM+\allowbreak{}stext & 0.07320 & 0.954 & 399.4 & 399.4 & 75946 \\
cls·B2 & DOM+\allowbreak{}sprompt & 0.08453 & 1.101 & 396.4 & 396.4 & 87948 \\
cls·B2 & SoM-image & 0.08456 & 1.102 & 411.0 & 411.0 & 87931 \\
red·B0 & DOM & 0.10147 & 1.000 & 572.0 & 552.5 & 93303 \\
red·B0 & SoM & 0.11045 & 1.089 & 461.4 & 451.6 & 103031 \\
red·B0 & Vision & 0.09807 & 0.966 & 449.8 & 418.5 & 88872 \\
red·B0 & DOM+\allowbreak{}stext & 0.10577 & 1.042 & 631.4 & 562.1 & 96214 \\
red·B0 & DOM+\allowbreak{}sprompt & 0.10163 & 1.002 & 498.4 & 447.7 & 93817 \\
red·B0 & SoM-image & 0.10814 & 1.066 & 562.0 & 532.0 & 99005 \\
red·B1 & DOM & 0.07330 & 1.000 & 602.6 & 602.6 & 76462 \\
red·B1 & SoM & 0.08000 & 1.091 & 609.1 & 609.1 & 83517 \\
red·B1 & Vision & 0.05240 & 0.715 & 456.6 & 456.6 & 53994 \\
red·B1 & DOM+\allowbreak{}stext & 0.06948 & 0.948 & 598.6 & 598.6 & 72222 \\
red·B1 & DOM+\allowbreak{}sprompt & 0.07656 & 1.044 & 614.0 & 614.0 & 79896 \\
red·B1 & SoM-image & 0.07480 & 1.020 & 616.6 & 616.6 & 77722 \\
red·B2 & DOM & 0.09479 & 1.000 & 669.9 & 669.9 & 99015 \\
red·B2 & SoM & 0.11160 & 1.177 & 623.0 & 623.0 & 117379 \\
red·B2 & Vision & 0.06833 & 0.721 & 550.0 & 550.0 & 70826 \\
red·B2 & DOM+\allowbreak{}stext & 0.08852 & 0.934 & 677.6 & 677.6 & 92346 \\
red·B2 & DOM+\allowbreak{}sprompt & 0.09940 & 1.049 & 599.3 & 599.3 & 104021 \\
red·B2 & SoM-image & 0.09451 & 0.997 & 640.0 & 640.0 & 98699 \\
WA·B0 & DOM & 0.07531 & 1.000 & 284.1 & 282.7 & 68616 \\
WA·B0 & SoM & 0.09110 & 1.210 & 272.1 & 271.7 & 84854 \\
WA·B0 & Vision & 0.08640 & 1.147 & 337.9 & 337.5 & 77786 \\
WA·B0 & DOM+\allowbreak{}stext & 0.08478 & 1.126 & 330.3 & 328.4 & 76591 \\
WA·B0 & DOM+\allowbreak{}sprompt & 0.07747 & 1.029 & 262.4 & 260.9 & 71176 \\
WA·B0 & SoM-image & 0.08498 & 1.128 & 324.0 & 320.9 & 77689 \\
WA·B1 & DOM & 0.06579 & 1.000 & 485.2 & 485.2 & 68491 \\
WA·B1 & SoM & 0.07944 & 1.208 & 494.8 & 494.8 & 82940 \\
WA·B1 & Vision & 0.04468 & 0.679 & 490.7 & 490.7 & 45878 \\
WA·B1 & DOM+\allowbreak{}stext & 0.06151 & 0.935 & 509.9 & 509.9 & 63727 \\
WA·B1 & DOM+\allowbreak{}sprompt & 0.07386 & 1.123 & 506.3 & 506.3 & 76839 \\
WA·B1 & SoM-image & 0.06659 & 1.012 & 510.6 & 510.6 & 68819 \\
\bottomrule
\end{tabular}
\end{table}


\section{Replicates and run-to-run noise}\label{app:noise}

\begin{figure}[tb]
\centering
\includegraphics[width=0.66\linewidth]{fig_ceilings}
\caption{The two ceilings, per cell (full table: Appendix Table~\ref{tab:t42}). Left: the six-run union ceiling against the best single mode; the tick is what one genuinely different arm adds, and the dark segment is what one rerun of an arm already in hand adds, where measured. Most of the apparent headroom is not attributable to representation diversity. Right: the cost ceiling at matched success, the bound that survives, because it adds no arm.}
\label{fig:ceilings}
\end{figure}


% Table source-order #27: Per-task label instability. Per-task label instability on \texttt{cls\
\begin{table}[tp]
\centering
\caption{Per-task label instability. Per-task label instability on \texttt{cls\_B0} (n=224): 49 tasks change outcome between two runs of the same condition. The rows a which-mode router could learn from are exactly the contested ones, and they carry almost all of the instability. \textbf{Caution: this table is computed on 2 replicated arms of one cell (B0.cls.dom, B0.cls.vision), rerun once.} The project's whole replicate inventory is three arms of that cell (the SoM pair landed after this product was generated and is not folded in here) plus five modes of \texttt{wa\_red\_B1} on a registered ten-task draw; no VWA-reddit cell and no B2 cell carries one at all, so every stability figure elsewhere is imported rather than measured. The headline enrichment has two defensible definitions and neither may be quoted alone: 17.4x defined over all 6 arms is correct for the claim (a router chooses among 6) but the flips are produced by rerunning 2 of them, so the same arms decide both membership and outcome; rebuilding the difficulty proxy from the other 4 breaks that circle and gives 3.95x. Source: \texttt{label\_instability.json}.}
\label{tab:t27}
\footnotesize
\setlength{\tabcolsep}{3pt}
\begin{tabular}{@{}
  >{\raggedright\arraybackslash}p{(\linewidth - 12\tabcolsep) * \real{0.1429}}
  >{\raggedright\arraybackslash}p{(\linewidth - 12\tabcolsep) * \real{0.1429}}
  >{\raggedright\arraybackslash}p{(\linewidth - 12\tabcolsep) * \real{0.1429}}
  >{\raggedright\arraybackslash}p{(\linewidth - 12\tabcolsep) * \real{0.1429}}
  >{\raggedright\arraybackslash}p{(\linewidth - 12\tabcolsep) * \real{0.1429}}
  >{\raggedright\arraybackslash}p{(\linewidth - 12\tabcolsep) * \real{0.1429}}
  >{\raggedright\arraybackslash}p{(\linewidth - 12\tabcolsep) * \real{0.1429}}@{}}
\toprule
\begin{minipage}[b]{\linewidth}\raggedright
stratum
\end{minipage} & \begin{minipage}[b]{\linewidth}\raggedright
tasks
\end{minipage} & \begin{minipage}[b]{\linewidth}\raggedright
share of cell
\end{minipage} & \begin{minipage}[b]{\linewidth}\raggedright
flipped
\end{minipage} & \begin{minipage}[b]{\linewidth}\raggedright
flip rate
\end{minipage} & \begin{minipage}[b]{\linewidth}\raggedright
share of all flips
\end{minipage} & \begin{minipage}[b]{\linewidth}\raggedright
enrichment vs complement
\end{minipage} \\
\midrule
which-mode label rows (any mode solved) & 97 & 43.3\% & 47 & 48.45\% & 95.9\% & \textbf{16.47x} \\
\ldots of those, the arms DISAGREE (the choice matters) & 88 & 39.3\% & 45 & 51.14\% & 91.8\% & \textbf{17.39x} \\
three-way channel decision is contested & 74 & 33.0\% & 36 & 48.65\% & 73.5\% & \textbf{16.54x} \\
exactly one mode solved it (label unambiguous) & 29 & 12.9\% & 15 & 51.72\% & 30.6\% & \textbf{17.59x} \\
complement: no mode solved, or all did & 136 & 60.7\% & 4 & 2.94\% & 8.2\% & \textbf{1.00x} \\
whole cell & 224 & 100.0\% & 49 & 21.88\% & 100.0\% & \textbf{7.44x} \\
\bottomrule
\end{tabular}
\end{table}

% Table source-order #14: Behavioural metrics against run-to-run noise. Behavioural metrics agai
\begin{table}[tp]
\centering
\caption{Behavioural metrics against run-to-run noise. Behavioural metrics against run-to-run movement, \texttt{B0\ x\ classifieds}, three replicated arms (dom, vision, som). \texttt{rerun\ band} is the largest \textbar metric(run A) - metric(run B)\textbar{} over those arms; \texttt{cross-mode\ spread} is max-min over the six modes. 22 of 25 metrics exceed the band, several by 5-22x. The exceptions are \texttt{mean\_latency\_canonical\_s} (0.84x), \texttt{mean\_latency\_s} (0.87x), \texttt{parse\_fail\_rate} (1.00x): i.e.~both latency metrics, which \texttt{latency\_decomposition} reaches independently by decomposing the step into model and container. Every other efficiency and behavioural claim in this paper is judged against a rerun band; these 26 metrics were not, until this table. One cell, one rerun per arm: a point estimate, not a threshold. Source: \texttt{replicate\_metric\_noise.json}.}
\label{tab:t14}
\footnotesize
\setlength{\tabcolsep}{3pt}
\begin{tabular}{@{}
  >{\raggedright\arraybackslash}p{(\linewidth - 10\tabcolsep) * \real{0.1667}}
  >{\raggedright\arraybackslash}p{(\linewidth - 10\tabcolsep) * \real{0.1667}}
  >{\raggedright\arraybackslash}p{(\linewidth - 10\tabcolsep) * \real{0.1667}}
  >{\raggedright\arraybackslash}p{(\linewidth - 10\tabcolsep) * \real{0.1667}}
  >{\raggedright\arraybackslash}p{(\linewidth - 10\tabcolsep) * \real{0.1667}}
  >{\raggedright\arraybackslash}p{(\linewidth - 10\tabcolsep) * \real{0.1667}}@{}}
\toprule
\begin{minipage}[b]{\linewidth}\raggedright
dimension
\end{minipage} & \begin{minipage}[b]{\linewidth}\raggedright
metric
\end{minipage} & \begin{minipage}[b]{\linewidth}\raggedright
cross-mode spread
\end{minipage} & \begin{minipage}[b]{\linewidth}\raggedright
rerun band
\end{minipage} & \begin{minipage}[b]{\linewidth}\raggedright
ratio
\end{minipage} & \begin{minipage}[b]{\linewidth}\raggedright
\textgreater{} a rerun?
\end{minipage} \\
\midrule
Outcome & \texttt{sr\_\allowbreak{}pct} & 11.607 & 2.232 & 5.20x & \textbf{yes} \\
Outcome & \texttt{n\_\allowbreak{}success} & 26.000 & 5.000 & 5.20x & \textbf{yes} \\
Outcome & \texttt{n\_\allowbreak{}unique\_\allowbreak{}solves} & --- & --- & --- & \emph{cross-mode by construction} \\
Macro & \texttt{n\_\allowbreak{}steps} & 2.567 & 0.469 & 5.48x & \textbf{yes} \\
Macro & \texttt{cap\_\allowbreak{}hit\_\allowbreak{}rate} & 0.062 & 0.045 & 1.40x & \textbf{yes} \\
Macro & \texttt{click\_\allowbreak{}frac} & 0.041 & 0.014 & 2.87x & \textbf{yes} \\
Macro & \texttt{type\_\allowbreak{}frac} & 0.078 & 0.019 & 4.21x & \textbf{yes} \\
Macro & \texttt{scroll\_\allowbreak{}frac} & 0.106 & 0.006 & 16.85x & \textbf{yes} \\
Macro & \texttt{search\_\allowbreak{}loop\_\allowbreak{}rate} & 0.121 & 0.018 & 6.75x & \textbf{yes} \\
Macro & \texttt{url\_\allowbreak{}revisit\_\allowbreak{}rate} & 0.081 & 0.022 & 3.66x & \textbf{yes} \\
Micro & \texttt{parse\_\allowbreak{}fail\_\allowbreak{}rate} & 0.002 & 0.002 & 1.00x & no \\
Micro & \texttt{action\_\allowbreak{}fail\_\allowbreak{}rate} & 0.068 & 0.014 & 4.87x & \textbf{yes} \\
Micro & \texttt{click\_\allowbreak{}fail\_\allowbreak{}rate} & 0.063 & 0.025 & 2.48x & \textbf{yes} \\
Micro & \texttt{type\_\allowbreak{}fail\_\allowbreak{}rate} & 0.030 & 0.006 & 5.28x & \textbf{yes} \\
Micro & \texttt{no\_\allowbreak{}change\_\allowbreak{}rate} & 0.053 & 0.009 & 5.82x & \textbf{yes} \\
Micro & \texttt{scroll\_\allowbreak{}inert\_\allowbreak{}rate} & 0.078 & 0.017 & 4.68x & \textbf{yes} \\
Micro & \texttt{noop\_\allowbreak{}inert\_\allowbreak{}rate} & 0.026 & 0.007 & 3.92x & \textbf{yes} \\
Micro & \texttt{visibility\_\allowbreak{}gap\_\allowbreak{}rate} & 0.016 & 0.011 & 1.56x & \textbf{yes} \\
Micro & \texttt{locator\_\allowbreak{}fallback\_\allowbreak{}rate} & 0.076 & 0.006 & 11.81x & \textbf{yes} \\
Micro & \texttt{action\_\allowbreak{}repeat\_\allowbreak{}frac} & 0.096 & 0.004 & 21.88x & \textbf{yes} \\
Micro & \texttt{finish\_\allowbreak{}rate} & 0.054 & 0.036 & 1.50x & \textbf{yes} \\
Efficiency & \texttt{mean\_\allowbreak{}cost\_\allowbreak{}usd} & 0.008 & 0.002 & 3.24x & \textbf{yes} \\
Efficiency & \texttt{cost\_\allowbreak{}rel\_\allowbreak{}dom} & 0.108 & 0.026 & 4.22x & \textbf{yes} \\
Efficiency & \texttt{mean\_\allowbreak{}latency\_\allowbreak{}s} & 19.562 & 22.488 & 0.87x & no \\
Efficiency & \texttt{mean\_\allowbreak{}latency\_\allowbreak{}canonical\_\allowbreak{}s} & 19.204 & 22.846 & 0.84x & no \\
Efficiency & \texttt{mean\_\allowbreak{}tokens} & 8697.978 & 2186.054 & 3.98x & \textbf{yes} \\
\bottomrule
\end{tabular}
\end{table}


\section{Efficiency and its estimands}\label{app:efficiency}

Two tables carry the efficiency point. Table~\ref{tab:t11} decomposes a latency figure: the model call is 22--67\% of a step, the rest is browser and container, and stripping the environment share changes which mode is fastest in 4 of 8 cells. Table~\ref{tab:t15} switches the denominator from per-attempt to per-success and the cheapest mode changes in 4 of the 6 cells with enough successes to divide by, with every pairwise interval overlapping. Both support the same statement: an efficiency claim without its estimand is ambiguous, and the estimand is usually left implicit.

% Table source-order #11: What a latency number contains. What a latency number contains. Every 
\begin{table}[tp]
\centering
\caption{What a latency number contains. What a latency number contains. Every latency figure elsewhere in this paper is the whole step; \texttt{backend\_infer} isolates the model call and was read by no analysis script until 2026-08-03. The model is 22--67\% of the measured time; the rest is the browser and the container, which \texttt{offsite\_navigation\_audit} measures at 1.69× between the two sites. Removing it changes which mode is fastest in 4 of 8 cells, and not at random: 4 of 5 reddit-family cells flip against 0 of 3 classifieds cells: the flips land where the container is slowest. A sentence naming the fastest mode is therefore partly a sentence about this deployment. What survives estimand choice is only that the two orderings disagree. Source: \texttt{latency\_decomposition.json}.}
\label{tab:t11}
\footnotesize
\setlength{\tabcolsep}{3pt}
\begin{tabular}{@{}
  >{\raggedright\arraybackslash}p{(\linewidth - 12\tabcolsep) * \real{0.1429}}
  >{\raggedright\arraybackslash}p{(\linewidth - 12\tabcolsep) * \real{0.1429}}
  >{\raggedright\arraybackslash}p{(\linewidth - 12\tabcolsep) * \real{0.1429}}
  >{\raggedright\arraybackslash}p{(\linewidth - 12\tabcolsep) * \real{0.1429}}
  >{\raggedright\arraybackslash}p{(\linewidth - 12\tabcolsep) * \real{0.1429}}
  >{\raggedright\arraybackslash}p{(\linewidth - 12\tabcolsep) * \real{0.1429}}
  >{\raggedright\arraybackslash}p{(\linewidth - 12\tabcolsep) * \real{0.1429}}@{}}
\toprule
\begin{minipage}[b]{\linewidth}\raggedright
cell
\end{minipage} & \begin{minipage}[b]{\linewidth}\raggedright
mean step (ms)
\end{minipage} & \begin{minipage}[b]{\linewidth}\raggedright
model call (ms)
\end{minipage} & \begin{minipage}[b]{\linewidth}\raggedright
model share
\end{minipage} & \begin{minipage}[b]{\linewidth}\raggedright
fastest by total
\end{minipage} & \begin{minipage}[b]{\linewidth}\raggedright
fastest by model only
\end{minipage} & \begin{minipage}[b]{\linewidth}\raggedright
same?
\end{minipage} \\
\midrule
B0·classifieds & 7,622 & 2,140 & \textbf{28.1\%} & DOM+\allowbreak{}sprompt & DOM+\allowbreak{}sprompt & yes \\
B0·reddit & 24,601 & 5,603 & \textbf{22.8\%} & Vision & SoM-image & \textbf{no} \\
B0·wa\_\allowbreak{}reddit & 16,112 & 3,551 & \textbf{22.0\%} & Vision & DOM+\allowbreak{}sprompt & \textbf{no} \\
B1·classifieds & 14,180 & 8,213 & \textbf{57.9\%} & Vision & Vision & yes \\
B1·reddit & 24,376 & 7,916 & \textbf{32.5\%} & Vision & DOM+\allowbreak{}stext & \textbf{no} \\
B1·wa\_\allowbreak{}reddit & 21,221 & 7,747 & \textbf{36.5\%} & DOM+\allowbreak{}sprompt & Vision & \textbf{no} \\
B2·classifieds & 14,739 & 9,908 & \textbf{67.2\%} & DOM+\allowbreak{}sprompt & DOM+\allowbreak{}sprompt & yes \\
B2·reddit & 22,863 & 9,131 & \textbf{39.9\%} & Vision & Vision & yes \\
\bottomrule
\end{tabular}
\end{table}

% Table source-order #15: Per-attempt versus per-success. Per-attempt versus per-success denomin
\begin{table}[tp]
\centering
\caption{Per-attempt versus per-success. Per-attempt versus per-success denominators. \texttt{content?\ =\ no} marks cells whose best mode has fewer than 10 successes: their ratios are directions at best and both B2 cells are in that state. Among the 6 cells with content the cheapest mode changes under the success denominator in 4 and the fastest in 3. \textbf{Caution:} Every pairwise CI overlaps, so this supports the methodological point that the denominator must be declared, not any ranking. Source: \texttt{outcome\_efficiency.json}.}
\label{tab:t15}
\footnotesize
\setlength{\tabcolsep}{3pt}
\begin{tabular}{@{}
  >{\raggedright\arraybackslash}p{(\linewidth - 12\tabcolsep) * \real{0.1429}}
  >{\raggedright\arraybackslash}p{(\linewidth - 12\tabcolsep) * \real{0.1429}}
  >{\raggedright\arraybackslash}p{(\linewidth - 12\tabcolsep) * \real{0.1429}}
  >{\raggedright\arraybackslash}p{(\linewidth - 12\tabcolsep) * \real{0.1429}}
  >{\raggedright\arraybackslash}p{(\linewidth - 12\tabcolsep) * \real{0.1429}}
  >{\raggedright\arraybackslash}p{(\linewidth - 12\tabcolsep) * \real{0.1429}}
  >{\raggedright\arraybackslash}p{(\linewidth - 12\tabcolsep) * \real{0.1429}}@{}}
\toprule
\begin{minipage}[b]{\linewidth}\raggedright
cell
\end{minipage} & \begin{minipage}[b]{\linewidth}\raggedright
content?
\end{minipage} & \begin{minipage}[b]{\linewidth}\raggedright
cheapest/attempt
\end{minipage} & \begin{minipage}[b]{\linewidth}\raggedright
cheapest/success
\end{minipage} & \begin{minipage}[b]{\linewidth}\raggedright
fastest/attempt
\end{minipage} & \begin{minipage}[b]{\linewidth}\raggedright
fastest/success
\end{minipage} & \begin{minipage}[b]{\linewidth}\raggedright
max solves
\end{minipage} \\
\midrule
cls·B0 & yes & Vision & Vision & SoM & SoM & 61 \\
red·B0 & yes & Vision & DOM & Vision & SoM & 30 \\
cls·B1 & yes & Vision & Vision & SoM & SoM & 32 \\
red·B1 & yes & Vision & SoM & Vision & SoM & 15 \\
cls·B2 & \textbf{no} & Vision & Vision & SoM & SoM & 5 \\
red·B2 & \textbf{no} & Vision & DOM & Vision & DOM & 8 \\
WA·B1 & yes & Vision & DOM+\allowbreak{}stext & DOM & DOM & 17 \\
WA·B0 & yes & DOM & DOM+\allowbreak{}stext & DOM+\allowbreak{}sprompt & DOM+\allowbreak{}stext & 37 \\
\bottomrule
\end{tabular}
\end{table}


\section{Routing, triage, and cascade}\label{app:routing}

\begin{figure}[tb]
\centering
\includegraphics[width=0.62\linewidth]{fig_partition_forest}
\caption{The zero-token partition (Appendix Table~\ref{tab:t17}): the screenshot's worth (Vision $-$ DOM) on regex-flagged tasks versus the rest. It is large and significant only on the capable classifieds backbones, and null on reddit. The predicate flags 71 of 224 classifieds tasks and 63 of 203 reddit tasks.}
\label{fig:partition}
\end{figure}


Table~\ref{tab:pb05} is the triage-label predictability that \S\ref{sec:lowerbound} builds on. Table~\ref{tab:t44} asks whether the two obstructions that section reports, too few labels and too few contested tasks, are two walls or one: both sets are subsets of the tasks something solves, and across the eight cells they track the best single mode's success rate at $\rho = 0.952$. Part of that association is structural rather than empirical, and the caption separates the two; what it licenses is the framing of \S\ref{sec:gap}, not a new effect.

% Paper B table #05: \caption{Predictability of the triage label. AUROC is cross-validated; the compa
\begin{table}[tp]
\centering
\caption{Predictability of the triage label. AUROC is cross-validated; the comparison column is the best single covariate used alone, so the margin isolates what the multivariate model adds. The reddit · B2 row is below chance and is the subject of \S\ref{sec:lowerbound}.}
\label{tab:pb05}
\footnotesize
\setlength{\tabcolsep}{3pt}
\begin{tabular}{@{}
  >{\raggedright\arraybackslash}p{(\linewidth - 10\tabcolsep) * \real{0.1667}}
  >{\raggedright\arraybackslash}p{(\linewidth - 10\tabcolsep) * \real{0.1667}}
  >{\raggedright\arraybackslash}p{(\linewidth - 10\tabcolsep) * \real{0.1667}}
  >{\raggedright\arraybackslash}p{(\linewidth - 10\tabcolsep) * \real{0.1667}}
  >{\raggedright\arraybackslash}p{(\linewidth - 10\tabcolsep) * \real{0.1667}}
  >{\raggedright\arraybackslash}p{(\linewidth - 10\tabcolsep) * \real{0.1667}}@{}}
\toprule
\begin{minipage}[b]{\linewidth}\raggedright
cell
\end{minipage} & \begin{minipage}[b]{\linewidth}\raggedright
tasks
\end{minipage} & \begin{minipage}[b]{\linewidth}\raggedright
solvable
\end{minipage} & \begin{minipage}[b]{\linewidth}\raggedright
AUROC (logistic)
\end{minipage} & \begin{minipage}[b]{\linewidth}\raggedright
best single covariate
\end{minipage} & \begin{minipage}[b]{\linewidth}\raggedright
margin
\end{minipage} \\
\midrule
classifieds · B0 & 224 & 43.3\% & 0.676 & 0.607 & +0.069 \\
reddit · B0 & 203 & 26.1\% & 0.666 & 0.612 & +0.054 \\
classifieds · B1 & 224 & 24.6\% & 0.717 & 0.627 & +0.090 \\
reddit · B1 & 203 & 11.8\% & 0.685 & 0.637 & +0.048 \\
classifieds · B2 & 224 & 7.1\% & 0.651 & 0.655 & −0.005 \\
reddit · B2 & 203 & 7.4\% & 0.483 & 0.711 & −0.228 \\
\bottomrule
\end{tabular}
\end{table}

% Table source-order #44: generated by scripts/analysis/export_ablation_tables.py
% Do not hand-edit the body; rerun the producer. Break-point fixes are the exception.
\begin{table}[tp]
\centering
\caption{Label supply and routing value are one quantity. Label supply and routing value are the same set. The rows a which-mode router can be trained on exist only where something succeeded; the rows where a per-task choice exists at all are those with more than one solver. Both are subsets of the solvable set, and across all 8 cells they track the best single mode's success rate with Spearman rho = 0.952 (exact permutation p = 0.0011), mean \textbar best SR - routable share\textbar{} = 1.65pp. So the two obstructions the lower bound reports are not independent walls but one wall, whose height is set by how much the agent can do. \textbf{Caution: part of this association is structural, not empirical}: both columns are functionals of one solve matrix. The non-structural content is the near-linearity: under mode independence the routable share would grow near-quadratically in the per-mode rate at these success levels, and it does not, which says task difficulty dominates mode-task matching. \textbf{Caution: the ratio is not constant}: 0.53-0.71 in the 6 cells above the floor, but half that (\texttt{red\_B2}, \texttt{cls\_B2}) on 3 and 4 routable tasks, where it is one task wide. \textbf{Caution: the cells are not independent}: they share three sites and three backbones, so the permutation null asks whether the pairing is arbitrary, not whether eight systems were sampled. Source: \texttt{supply\_value\_coupling.json}.}
\label{tab:t44}
\footnotesize
\setlength{\tabcolsep}{3pt}
\begin{tabular}{@{}
  >{\raggedright\arraybackslash}p{(\linewidth - 12\tabcolsep) * \real{0.1429}}
  >{\raggedright\arraybackslash}p{(\linewidth - 12\tabcolsep) * \real{0.1429}}
  >{\raggedright\arraybackslash}p{(\linewidth - 12\tabcolsep) * \real{0.1429}}
  >{\raggedright\arraybackslash}p{(\linewidth - 12\tabcolsep) * \real{0.1429}}
  >{\raggedright\arraybackslash}p{(\linewidth - 12\tabcolsep) * \real{0.1429}}
  >{\raggedright\arraybackslash}p{(\linewidth - 12\tabcolsep) * \real{0.1429}}
  >{\raggedright\arraybackslash}p{(\linewidth - 12\tabcolsep) * \real{0.1429}}@{}}
\toprule
\begin{minipage}[b]{\linewidth}\raggedright
cell
\end{minipage} & \begin{minipage}[b]{\linewidth}\raggedright
n
\end{minipage} & \begin{minipage}[b]{\linewidth}\raggedright
best single mode
\end{minipage} & \begin{minipage}[b]{\linewidth}\raggedright
best SR
\end{minipage} & \begin{minipage}[b]{\linewidth}\raggedright
routable set (\textgreater1 solver)
\end{minipage} & \begin{minipage}[b]{\linewidth}\raggedright
solvable (any mode)
\end{minipage} & \begin{minipage}[b]{\linewidth}\raggedright
routable / solvable
\end{minipage} \\
\midrule
\texttt{wa\_red\_B0} & 104 & DOM+stext & 35.58\% & 34.6\% (36) & 51.9\% & 0.67 \\
\texttt{cls\_B0} & 224 & SoM & 27.23\% & 30.4\% (68) & 43.3\% & 0.70 \\
\texttt{wa\_red\_B1} & 104 & DOM & 16.35\% & 17.3\% (18) & 30.8\% & 0.56 \\
\texttt{red\_B0} & 203 & SoM & 14.78\% & 17.7\% (36) & 26.1\% & 0.68 \\
\texttt{cls\_B1} & 224 & SoM & 14.29\% & 12.9\% (29) & 24.6\% & 0.53 \\
\texttt{red\_B1} & 203 & SoM & 7.39\% & 8.4\% (17) & 11.8\% & 0.71 \\
\texttt{red\_B2} & 203 & DOM & 3.94\% & 1.5\% (3) & 7.4\% & 0.20 \\
\texttt{cls\_B2} & 224 & SoM & 2.23\% & 1.8\% (4) & 7.1\% & 0.25 \\
\bottomrule
\end{tabular}
\end{table}



% Trailing sections of this appendix live in sections/appendix_pb.tex.

Table~\ref{tab:t19} turns the ex-ante partition (Table~\ref{tab:t17}) into a policy. Table~\ref{tab:t20}'s cascade beats always-rich at no operating point in any comparable cell, and Table~\ref{tab:t21} is the table that says what that means: against the \emph{same budget spent at random} the confidence ranking wins nearly everywhere, so the signal is informative and still not enough. Table~\ref{tab:t22} adds the benchmark's own difficulty annotation for a mean $\Delta$AUROC of +0.0024. Table~\ref{tab:t23} shows the feature a practitioner would reach for first looks like it should help and does not. Table~\ref{tab:t33} pools the backbones and routes by cost tier instead of by task.

% Resolved 2026-08-05: the two papers merged, so the former companion-overlap
% cut-candidates (t20-t22, t33) are retained as part of the five-route lower
% bound; t33's truncated caption and damaged cells were rewritten 4a605ab.

The bound the rest of this section does not reach is Table~\ref{tab:t42} in \S\ref{sec:upperbound}. The oracle stated against the best single mode, and its triage/route decomposition, are below, followed by the rerun comparator (Table~\ref{tab:t18}), the per-cell fusion-premium intervals behind Figure~\ref{fig:fusion} (Table~\ref{tab:t16}), the which-mode label supply (Table~\ref{tab:pb03}), and the per-cell partition intervals behind Figure~\ref{fig:partition} (Table~\ref{tab:t17}).

% Table source-order #42: What a perfect per-task choice could buy.
\begin{table}[tbp]
\centering
\caption{What a perfect per-task choice could buy. Two ceilings, and only one of them survives its own control. \emph{Ceiling: any mode solves it} is what a perfect per-task choice could reach; it runs 7.1-51.9\% against a best single mode of 2.2-35.6\%. \textbf{Caution:} That column is a six-arm union against a one-arm baseline. The arm-matched comparison is the last two columns: adding the single best distinct arm buys +1.97 to +7.14pp, and rerunning an arm already in hand buys a draw in the same range wherever a replicate exists -- so the headroom cannot be attributed to representation diversity rather than to resampling. The union of \emph{five} reruns has never been measured, so the split at higher arm counts is unknown, not estimated. \emph{Same tasks, lower cost} keeps the best mode everywhere and sends only the tasks no mode solves to the cheapest one: success is unchanged by construction and cost falls 9.5-30.6\% in 8 of 8 cells. That ceiling is immune to the arm-count objection because it adds no arms. Why both are hard to reach: no mode solves 48.1-92.9\% of tasks, and the set where a per-task choice even exists is 1.5-34.6\% of the cell. Leaked successes are kept here, as in every other table; 6 scored successes on reddit were credited without the episode ever visiting the forum the evaluator reads, and zeroing them (denominator unchanged) is reported as a sensitivity analysis rather than folded in: the detection criterion is a lower bound, so it cannot license a corrected point estimate. Both sets of figures are in the product. Source: \texttt{routing\_ceiling.json}}
\label{tab:t42}
\footnotesize
\setlength{\tabcolsep}{4pt}
\begin{tabular}{@{}
  >{\raggedright\arraybackslash}p{(\linewidth - 14\tabcolsep) * \real{0.1250}}
  >{\raggedright\arraybackslash}p{(\linewidth - 14\tabcolsep) * \real{0.1250}}
  >{\raggedright\arraybackslash}p{(\linewidth - 14\tabcolsep) * \real{0.1250}}
  >{\raggedright\arraybackslash}p{(\linewidth - 14\tabcolsep) * \real{0.1250}}
  >{\raggedright\arraybackslash}p{(\linewidth - 14\tabcolsep) * \real{0.1250}}
  >{\raggedright\arraybackslash}p{(\linewidth - 14\tabcolsep) * \real{0.1250}}
  >{\raggedright\arraybackslash}p{(\linewidth - 14\tabcolsep) * \real{0.1250}}
  >{\raggedright\arraybackslash}p{(\linewidth - 14\tabcolsep) * \real{0.1250}}@{}}
\toprule
\begin{minipage}[b]{\linewidth}\raggedright
cell
\end{minipage} & \begin{minipage}[b]{\linewidth}\raggedright
n
\end{minipage} & \begin{minipage}[b]{\linewidth}\raggedright
best single mode
\end{minipage} & \begin{minipage}[b]{\linewidth}\raggedright
ceiling: any mode solves
\end{minipage} & \begin{minipage}[b]{\linewidth}\raggedright
headroom
\end{minipage} & \begin{minipage}[b]{\linewidth}\raggedright
same tasks, lower cost
\end{minipage} & \begin{minipage}[b]{\linewidth}\raggedright
+1 arm
\end{minipage} & \begin{minipage}[b]{\linewidth}\raggedright
rerun once
\end{minipage} \\
\midrule
WA·B0 & 104 & DOM+\allowbreak{}stext 35.58\% & \textbf{51.92\%} & +16.35pp & \textbf{-10.7\%} & +5.77 & 2.00-4.00 \\
cls·B0 & 224 & SoM 27.23\% & \textbf{43.30\%} & +16.07pp & \textbf{-12.8\%} & +7.14 & 4.91-7.59 \\
WA·B1 & 104 & DOM 16.35\% & \textbf{30.77\%} & +14.42pp & \textbf{-26.8\%} & +4.81 & 2.00-4.00 \\
red·B0 & 203 & SoM 14.78\% & \textbf{26.11\%} & +11.33pp & \textbf{-9.5\%} & +4.93 & -- \\
cls·B1 & 224 & SoM 14.29\% & \textbf{24.55\%} & +10.27pp & \textbf{-19.4\%} & +4.91 & -- \\
red·B1 & 203 & SoM 7.39\% & \textbf{11.82\%} & +4.43pp & \textbf{-30.6\%} & +1.97 & -- \\
red·B2 & 203 & DOM 3.94\% & \textbf{7.39\%} & +3.45pp & \textbf{-26.4\%} & +1.97 & -- \\
cls·B2 & 224 & SoM 2.23\% & \textbf{7.14\%} & +4.91pp & \textbf{-21.3\%} & +2.23 & -- \\
\bottomrule
\end{tabular}
\end{table}

% Paper B table #01: \caption{The routing ceiling, over classifieds (cls) and reddit (red). The last
\begin{table}[tp]
\centering
\caption{The routing ceiling, over classifieds (cls) and reddit (red). The last three columns give success rate (\%) and mean per-episode billed cost (USD). The oracle picks the cheapest mode that solved each task, so its success rate equals the solvable column by construction; ``cheapest'' is the lowest-mean-cost mode applied everywhere. Cost is never compared across backbones.}
\label{tab:pb01}
\footnotesize
\setlength{\tabcolsep}{3pt}
\begin{tabular}{@{}
  >{\raggedright\arraybackslash}p{(\linewidth - 10\tabcolsep) * \real{0.1667}}
  >{\raggedright\arraybackslash}p{(\linewidth - 10\tabcolsep) * \real{0.1667}}
  >{\raggedright\arraybackslash}p{(\linewidth - 10\tabcolsep) * \real{0.1667}}
  >{\raggedright\arraybackslash}p{(\linewidth - 10\tabcolsep) * \real{0.1667}}
  >{\raggedright\arraybackslash}p{(\linewidth - 10\tabcolsep) * \real{0.1667}}
  >{\raggedright\arraybackslash}p{(\linewidth - 10\tabcolsep) * \real{0.1667}}@{}}
\toprule
\begin{minipage}[b]{\linewidth}\raggedright
cell
\end{minipage} & \begin{minipage}[b]{\linewidth}\raggedright
n
\end{minipage} & \begin{minipage}[b]{\linewidth}\raggedright
solvable
\end{minipage} & \begin{minipage}[b]{\linewidth}\raggedright
best single
\end{minipage} & \begin{minipage}[b]{\linewidth}\raggedright
cheapest
\end{minipage} & \begin{minipage}[b]{\linewidth}\raggedright
oracle
\end{minipage} \\
\midrule
cls · B0 & 224 & 43.3\% & 27.23 / 0.07236 & 25.00 / 0.06481 & 43.30 / 0.05777 \\
red · B0 & 203 & 26.1\% & 14.78 / 0.11045 & 7.39 / 0.09807 & 26.11 / 0.09534 \\
cls · B1 & 224 & 24.6\% & 14.29 / 0.06028 & 12.50 / 0.04316 & 24.55 / 0.04171 \\
red · B1 & 203 & 11.8\% & 7.39 / 0.08000 & 2.46 / 0.05240 & 11.82 / 0.05178 \\
cls · B2 & 224 & 7.1\% & 2.23 / 0.09075 & 2.23 / 0.07065 & 7.14 / 0.06953 \\
red · B2 & 203 & 7.4\% & 3.94 / 0.09479 & 1.97 / 0.06833 & 7.39 / 0.06958 \\
\bottomrule
\end{tabular}
\end{table}

% Paper B table #02: \caption{The ceiling decomposed. Both columns are deltas against the best single
\begin{table}[tp]
\centering
\caption{The ceiling decomposed. Both columns are deltas against the best single mode. The triage half is accuracy-neutral by construction and carries 9.5--30.6\% of cost saving; the route half carries 3.45--16.35pp of success rate and much less cost. The two WebArena rows were added 2026-08-04 and are the favourable end of both halves: \texttt{WA$\cdot$B0} carries the largest route-half headroom (+16.35pp) and the largest route-half cost saving (−16.2\%) anywhere in the study, so the six-cell version understated the ceiling on the cells where routing has most to work with. The two halves require different labels, and the rest of the paper fails on one obstruction per half.}
\label{tab:pb02}
\footnotesize
\setlength{\tabcolsep}{3pt}
\begin{tabular}{@{}lll@{}}
\toprule
cell & triage half: ΔSR / Δcost & route half: ΔSR / Δcost \\
\midrule
classifieds · B0 & ±0.00pp / −12.8\% & +16.07pp / −7.4\% \\
reddit · B0 & ±0.00pp / −9.5\% & +11.33pp / −4.2\% \\
classifieds · B1 & ±0.00pp / −19.4\% & +10.27pp / −11.4\% \\
reddit · B1 & ±0.00pp / −30.6\% & +4.43pp / −4.7\% \\
classifieds · B2 & ±0.00pp / −21.3\% & +4.91pp / −2.1\% \\
reddit · B2 & ±0.00pp / −26.4\% & +3.45pp / −0.2\% \\
WA reddit · B0 & ±0.00pp / −10.7\% & +16.35pp / −16.2\% \\
WA reddit · B1 & ±0.00pp / −26.8\% & +14.42pp / −6.7\% \\
\bottomrule
\end{tabular}
\end{table}

% Table source-order #18: New representation versus a rerun. Is a new representation worth more 
\begin{table}[tp]
\centering
\caption{New representation versus a rerun. Is a new representation worth more than a rerun? Both middle columns are the same functional at the same arm count (\texttt{\textbar{}\{added\}\ ∖\ \{baseline\}\textbar{}\ /\ n}) so they are directly comparable; only the \emph{source} of the extra arm differs. The band is 3 rerun pairs on one cell (\texttt{B0\ ×\ classifieds}, n=224), one each for DOM, SoM, Vision: so the rows without a band have no comparator at all, and the band itself is 3 draws rather than a bound. Since the SoM replicate landed 2026-08-03 the band is no longer extrapolated onto an unreplicated arm: both the fused mode this table's best-single column keeps selecting and the arm the comparison adds now carry their own measured floor, and adding the third pair left the band unmoved. Source: \texttt{noise\_floor\_inventory.json}}
\label{tab:t18}
\footnotesize
\setlength{\tabcolsep}{3pt}
\begin{tabular}{@{}
  >{\raggedright\arraybackslash}p{(\linewidth - 8\tabcolsep) * \real{0.2000}}
  >{\raggedright\arraybackslash}p{(\linewidth - 8\tabcolsep) * \real{0.2000}}
  >{\raggedright\arraybackslash}p{(\linewidth - 8\tabcolsep) * \real{0.2000}}
  >{\raggedright\arraybackslash}p{(\linewidth - 8\tabcolsep) * \real{0.2000}}
  >{\raggedright\arraybackslash}p{(\linewidth - 8\tabcolsep) * \real{0.2000}}@{}}
\toprule
\begin{minipage}[b]{\linewidth}\raggedright
cell
\end{minipage} & \begin{minipage}[b]{\linewidth}\raggedright
best single
\end{minipage} & \begin{minipage}[b]{\linewidth}\raggedright
+1 distinct arm
\end{minipage} & \begin{minipage}[b]{\linewidth}\raggedright
+1 rerun (measured floor)
\end{minipage} & \begin{minipage}[b]{\linewidth}\raggedright
verdict
\end{minipage} \\
\midrule
cls·B0 & SoM at 27.23 & +7.14 (DOM) & 4.91 -- 7.59pp & \textbf{inside the rerun band} \\
cls·B1 & SoM at 14.29 & +4.91 (Vision) & --- & no floor on this cell \\
cls·B2 & SoM at 2.23 & +2.23 (Vision) & --- & no floor on this cell \\
red·B0 & SoM at 14.78 & +4.93 (DOM) & --- & no floor on this cell \\
red·B1 & SoM at 7.39 & +1.97 (DOM+\allowbreak{}sprompt) & --- & no floor on this cell \\
red·B2 & DOM at 3.94 & +1.97 (Vision) & --- & no floor on this cell \\
WA·B0 & DOM+\allowbreak{}stext at 35.58 & +5.77 (DOM) & --- & no floor on this cell \\
WA·B1 & DOM at 16.35 & +4.81 (DOM+\allowbreak{}stext) & 0.00 -- 10.00pp \emph{(pooled would read 2.00--4.00)} & \textbf{inside the rerun band} \\
\bottomrule
\end{tabular}
\end{table}

\begin{figure}[tp]
\centering
\includegraphics[width=\linewidth]{fig_fusion_forest}
\caption{Fusion premium per cell (Appendix Table~\ref{tab:t16}), paired bootstrap 95\% CIs. The shaded region is the rerun threshold of \S\ref{sec:noise}: 3.8--4.2pp, the level a mean difference must reach before a single repetition of one arm would be unlikely to produce it. Against the workload-matched channel no cell clears it (the largest such premium is +2.23pp); SoM's decisive wins are all against the channel mismatched to that workload.}
\label{fig:fusion}
\end{figure}
% Table source-order #16: Fusion premium against the rerun band. Fusion premium (pp). Comparator
\begin{table}[tp]
\centering
\caption{Fusion premium against the rerun band. Fusion premium (pp). Comparators are fixed a priori, not per-cell maxima. Paired bootstrap over tasks, 10,000 resamples. Read against the measured rerun band 0.89--2.23pp, not against zero: a premium must beat what repetition delivers for the same money. No cell clears the band; \texttt{cls\_B0}'s +2.23 \emph{equals} its upper edge (both are 5/224). Source: \texttt{fusion\_premium.json}}
\label{tab:t16}
\footnotesize
\setlength{\tabcolsep}{3pt}
\begin{tabular}{@{}llllll@{}}
\toprule
cell & n & SoM − Vision & 95\% CI & SoM − DOM & 95\% CI \\
\midrule
cls·B0 & 224 & +2.23 & {[}-2.68, +7.59{]} & +9.82 & {[}+3.57, +16.07{]} \\
cls·B1 & 224 & +1.79 & {[}-2.68, +6.25{]} & +8.04 & {[}+3.57, +12.95{]} \\
cls·B2 & 224 & +0.00 & {[}-2.68, +2.68{]} & +0.89 & {[}-1.34, +3.57{]} \\
red·B0 & 203 & +7.39 & {[}+2.46, +12.32{]} & +0.49 & {[}-3.94, +4.93{]} \\
red·B1 & 203 & +4.93 & {[}+1.48, +8.87{]} & +1.48 & {[}-1.48, +4.43{]} \\
red·B2 & 203 & -0.99 & {[}-3.45, +1.48{]} & -2.96 & {[}-5.91, -0.49{]} \\
WA·B1 & 104 & +3.85 & {[}-1.92, +9.62{]} & -2.88 & {[}-9.62, +2.88{]} \\
WA·B0 & 104 & +2.88 & {[}-5.77, +11.54{]} & -4.81 & {[}-12.50, +2.88{]} \\
\bottomrule
\end{tabular}
\end{table}

% Paper B table #03: \caption{Which-mode label supply. A label exists only where some mode solved the
\begin{table}[tp]
\centering
\caption{Which-mode label supply. A label exists only where some mode solved the task, so the labelled column equals the number of solved tasks. Six classes have to be discriminated from between 15 and 97 examples.}
\label{tab:pb03}
\footnotesize
\setlength{\tabcolsep}{3pt}
\begin{tabular}{@{}lllll@{}}
\toprule
cell & scored tasks & labelled & solvable & classes present \\
\midrule
classifieds · B0 & 224 & 97 & 43.3\% & 6/6 \\
reddit · B0 & 203 & 53 & 26.1\% & 6/6 \\
classifieds · B1 & 224 & 55 & 24.6\% & 6/6 \\
reddit · B1 & 203 & 24 & 11.8\% & 6/6 \\
classifieds · B2 & 224 & 16 & 7.1\% & 5/6 \\
reddit · B2 & 203 & 15 & 7.4\% & 5/6 \\
\bottomrule
\end{tabular}
\end{table}

% Paper B table #04: \caption{Classes surviving a minimum of ten training rows per class in a five-fo
\begin{table}[tp]
\centering
\caption{Classes surviving a minimum of ten training rows per class in a five-fold split. Fewer than two surviving classes leaves nothing to discriminate between, which is the state of four of the six cells.}
\label{tab:pb04}
\footnotesize
\setlength{\tabcolsep}{3pt}
\begin{tabular}{@{}lllll@{}}
\toprule
cell & labels & classes & classes surviving the filter & trainable \\
\midrule
classifieds · B0 & 97 & 6 & 3 (DOM, DOM+\allowbreak{}sprompt, SoM) & yes \\
reddit · B0 & 53 & 6 & 1 (DOM) & no \\
classifieds · B1 & 55 & 6 & 2 (DOM, SoM) & yes \\
reddit · B1 & 24 & 6 & 0 & no \\
classifieds · B2 & 16 & 5 & 0 & no \\
reddit · B2 & 15 & 5 & 0 & no \\
\bottomrule
\end{tabular}
\end{table}

% Table source-order #17: Ex-ante visual-intent partition. Ex-ante partition. The predicate is a
\begin{table}[tp]
\centering
\caption{Ex-ante visual-intent partition. Ex-ante partition. The predicate is a regex over the task intent plus `carries no reference image': both read the task config, so it costs no tokens and needs no episode. On classifieds the screenshot is worth an order of magnitude more on the flagged tasks than on the rest, and the flagged/rest split is significant/not respectively on the two capable backbones. \textbf{Caution:} WebArena is omitted: the predicate fires on only 5 of 104 tasks there and none is solved by any mode, so the cells are degenerate rather than null. Source: \texttt{visual\_intent\_routing.json}.}
\label{tab:t17}
\footnotesize
\setlength{\tabcolsep}{3pt}
\begin{tabular}{@{}
  >{\raggedright\arraybackslash}p{(\linewidth - 12\tabcolsep) * \real{0.1429}}
  >{\raggedright\arraybackslash}p{(\linewidth - 12\tabcolsep) * \real{0.1429}}
  >{\raggedright\arraybackslash}p{(\linewidth - 12\tabcolsep) * \real{0.1429}}
  >{\raggedright\arraybackslash}p{(\linewidth - 12\tabcolsep) * \real{0.1429}}
  >{\raggedright\arraybackslash}p{(\linewidth - 12\tabcolsep) * \real{0.1429}}
  >{\raggedright\arraybackslash}p{(\linewidth - 12\tabcolsep) * \real{0.1429}}
  >{\raggedright\arraybackslash}p{(\linewidth - 12\tabcolsep) * \real{0.1429}}@{}}
\toprule
\begin{minipage}[b]{\linewidth}\raggedright
cell
\end{minipage} & \begin{minipage}[b]{\linewidth}\raggedright
flagged
\end{minipage} & \begin{minipage}[b]{\linewidth}\raggedright
arm
\end{minipage} & \begin{minipage}[b]{\linewidth}\raggedright
Δ vs DOM on flagged
\end{minipage} & \begin{minipage}[b]{\linewidth}\raggedright
95\% CI
\end{minipage} & \begin{minipage}[b]{\linewidth}\raggedright
Δ on the rest
\end{minipage} & \begin{minipage}[b]{\linewidth}\raggedright
95\% CI
\end{minipage} \\
\midrule
cls·B0 & 71 & vision & +22.54 & {[}+9.86, +33.80{]} & +0.65 & {[}-5.88, +7.84{]} \\
cls·B0 & 71 & som & +19.72 & {[}+7.04, +32.39{]} & +5.23 & {[}-1.31, +12.42{]} \\
cls·B1 & 71 & vision & +16.90 & {[}+8.45, +25.35{]} & +1.31 & {[}-3.92, +6.54{]} \\
cls·B1 & 71 & som & +12.68 & {[}+5.63, +21.13{]} & +5.88 & {[}+0.00, +11.76{]} \\
cls·B2 & 71 & vision & +1.41 & {[}-4.23, +7.04{]} & +0.65 & {[}-1.31, +3.27{]} \\
cls·B2 & 71 & som & +0.00 & {[}-5.63, +5.63{]} & +1.31 & {[}-1.31, +3.92{]} \\
red·B0 & 63 & vision & -3.17 & {[}-14.29, +7.94{]} & -8.57 & {[}-14.29, -2.86{]} \\
red·B0 & 63 & som & +0.00 & {[}-9.52, +9.52{]} & +0.71 & {[}-4.29, +5.71{]} \\
red·B1 & 63 & vision & -3.17 & {[}-9.52, +3.17{]} & -3.57 & {[}-7.86, +0.71{]} \\
red·B1 & 63 & som & +1.59 & {[}-3.17, +7.94{]} & +1.43 & {[}-2.14, +5.00{]} \\
red·B2 & 63 & vision & +4.76 & {[}-1.59, +12.70{]} & -5.00 & {[}-8.57, -1.43{]} \\
red·B2 & 63 & som & -1.59 & {[}-4.76, +0.00{]} & -3.57 & {[}-7.14, +0.00{]} \\
\bottomrule
\end{tabular}
\end{table}


% Table source-order #19: Routing policies on the 3-axis frontier. Is routing worth it? \texttt{
\begin{table}[tp]
\centering
\caption{Routing policies on the 3-axis frontier. Is routing worth it? \texttt{rule} policies send the ex-ante-flagged tasks to one arm and the rest to another; the partition is a regex over the task intent, so nothing is learned and there is no in-sample optimism. `On frontier' means nothing dominates it, not that it is preferable: on \texttt{cls·B0} all three rule policies sit between always-SoM and always-Vision, worse on every axis than one or the other. Cost/latency are per-attempt cell means, within-cell comparable only. Source: \texttt{rule\_routing\_pareto.json}.}
\label{tab:t19}
\footnotesize
\setlength{\tabcolsep}{3pt}
\begin{tabular}{@{}llllll@{}}
\toprule
cell & policy & SR & cost & latency & on frontier \\
\midrule
cls·B0 & always-SoM & 27.23 & 0.07236 & 106.0s & yes \\
cls·B0 & always-Vision & 25.00 & 0.06481 & 125.2s & yes \\
cls·B0 & rule: flag→Vision else DOM & 24.55 & 0.06809 & 117.6s & yes \\
cls·B0 & rule: flag→SoM else DOM+\allowbreak{}stext & 24.11 & 0.07019 & 115.8s & yes \\
cls·B0 & rule: flag→SoM else DOM & 23.66 & 0.07049 & 111.5s & yes \\
cls·B0 & always-DOM & 17.41 & 0.06962 & 114.1s & no \\
cls·B1 & always-SoM & 14.29 & 0.06028 & 262.0s & yes \\
cls·B1 & always-Vision & 12.50 & 0.04316 & 269.8s & yes \\
cls·B1 & rule: flag→Vision else DOM & 11.61 & 0.05433 & 296.5s & no \\
cls·B1 & rule: flag→SoM else DOM+\allowbreak{}stext & 11.16 & 0.05926 & 297.2s & no \\
cls·B1 & rule: flag→SoM else DOM & 10.27 & 0.05976 & 294.0s & no \\
cls·B1 & always-DOM & 6.25 & 0.05951 & 308.9s & no \\
cls·B2 & always-SoM & 2.23 & 0.09075 & 374.3s & yes \\
cls·B2 & always-Vision & 2.23 & 0.07065 & 417.8s & yes \\
cls·B2 & rule: flag→Vision else DOM & 1.79 & 0.07483 & 407.2s & yes \\
cls·B2 & always-DOM & 1.34 & 0.07676 & 402.3s & yes \\
cls·B2 & rule: flag→SoM else DOM & 1.34 & 0.08120 & 393.4s & no \\
cls·B2 & rule: flag→SoM else DOM+\allowbreak{}stext & 1.34 & 0.07876 & 391.5s & yes \\
red·B0 & always-SoM & 14.78 & 0.11045 & 451.6s & yes \\
red·B0 & always-DOM & 14.29 & 0.10147 & 552.5s & yes \\
red·B0 & rule: flag→SoM else DOM & 14.29 & 0.10426 & 521.2s & yes \\
red·B0 & rule: flag→Vision else DOM & 13.30 & 0.10041 & 510.9s & yes \\
red·B0 & rule: flag→SoM else DOM+\allowbreak{}stext & 13.30 & 0.10722 & 527.8s & no \\
red·B0 & always-Vision & 7.39 & 0.09807 & 418.5s & yes \\
red·B1 & always-SoM & 7.39 & 0.08000 & 609.1s & yes \\
red·B1 & rule: flag→SoM else DOM+\allowbreak{}stext & 6.90 & 0.07275 & 601.8s & yes \\
red·B1 & rule: flag→SoM else DOM & 6.40 & 0.07538 & 604.6s & no \\
red·B1 & always-DOM & 5.91 & 0.07330 & 602.6s & no \\
red·B1 & rule: flag→Vision else DOM & 4.93 & 0.06682 & 557.3s & yes \\
red·B1 & always-Vision & 2.46 & 0.05240 & 456.6s & yes \\
red·B2 & rule: flag→Vision else DOM & 5.42 & 0.08658 & 632.7s & yes \\
red·B2 & always-DOM & 3.94 & 0.09479 & 669.9s & no \\
red·B2 & rule: flag→SoM else DOM & 3.45 & 0.10001 & 655.4s & no \\
red·B2 & always-Vision & 1.97 & 0.06833 & 550.0s & yes \\
red·B2 & rule: flag→SoM else DOM+\allowbreak{}stext & 1.48 & 0.09568 & 660.7s & no \\
red·B2 & always-SoM & 0.99 & 0.11160 & 623.0s & no \\
\bottomrule
\end{tabular}
\end{table}

% Table source-order #20: Confidence-triggered cascade. Confidence-triggered cascade, vision → s
\begin{table}[tp]
\centering
\caption{Confidence-triggered cascade. Confidence-triggered cascade, vision → som. The escalation decision sees only the cheap run's own episode: no outcome, no rich-run information. In 0 of the 6 comparable cells does any operating point beat always-rich on success rate; the cells marked \emph{rich is worse} are excluded because the cascade's premise fails there, and that exclusion rule (\texttt{cheap\_sr\ \textgreater{}=\ rich\_sr}) is outcome-dependent. \textbf{Caution:} The two counts are split for a reason. A point also `wins' by matching always-rich's SR at lower cost, and every such point here is a tie artefact: on \texttt{WA·B1} all of them sit at exactly the rich arm's SR while 44-60 of 104 episodes are tied at the cutoff, so the ranking falls through to task id and the reachable SR spans 8.65-14.42. Counting the two together is what let this table print a WA win for months while the prose said the WA exception was withdrawn: the rows disagreed with the sentence and nobody could see it, because the WA rows were dropped before render until 2026-08-03. \textbf{Caution:} Every number is an offline splice: an escalated task takes its outcome from a standalone rich run, whereas a real cascade would start the rich episode after the cheap one had already acted on a stateful site. That sequential outcome is unobserved in this project. Source: \texttt{confidence\_cascade\_with\_wa.json}}
\label{tab:t20}
\footnotesize
\setlength{\tabcolsep}{3pt}
\begin{tabular}{@{}
  >{\raggedright\arraybackslash}p{(\linewidth - 16\tabcolsep) * \real{0.1111}}
  >{\raggedright\arraybackslash}p{(\linewidth - 16\tabcolsep) * \real{0.1111}}
  >{\raggedright\arraybackslash}p{(\linewidth - 16\tabcolsep) * \real{0.1111}}
  >{\raggedright\arraybackslash}p{(\linewidth - 16\tabcolsep) * \real{0.1111}}
  >{\raggedright\arraybackslash}p{(\linewidth - 16\tabcolsep) * \real{0.1111}}
  >{\raggedright\arraybackslash}p{(\linewidth - 16\tabcolsep) * \real{0.1111}}
  >{\raggedright\arraybackslash}p{(\linewidth - 16\tabcolsep) * \real{0.1111}}
  >{\raggedright\arraybackslash}p{(\linewidth - 16\tabcolsep) * \real{0.1111}}
  >{\raggedright\arraybackslash}p{(\linewidth - 16\tabcolsep) * \real{0.1111}}@{}}
\toprule
\begin{minipage}[b]{\linewidth}\raggedright
cell
\end{minipage} & \begin{minipage}[b]{\linewidth}\raggedright
n
\end{minipage} & \begin{minipage}[b]{\linewidth}\raggedright
cheap SR
\end{minipage} & \begin{minipage}[b]{\linewidth}\raggedright
rich SR
\end{minipage} & \begin{minipage}[b]{\linewidth}\raggedright
always-rich cost
\end{minipage} & \begin{minipage}[b]{\linewidth}\raggedright
oracle SR
\end{minipage} & \begin{minipage}[b]{\linewidth}\raggedright
comparable?
\end{minipage} & \begin{minipage}[b]{\linewidth}\raggedright
beats always-rich --- strictly / SR-tied
\end{minipage} & \begin{minipage}[b]{\linewidth}\raggedright
signals dropped
\end{minipage} \\
\midrule
cls·B0 & 224 & 25.00 & 27.23 & 1.116× & 33.93 & yes & 0 / 0 & 2 \\
cls·B1 & 224 & 12.50 & 14.29 & 1.397× & 19.20 & yes & 0 / 0 & 0 \\
cls·B2 & 224 & 2.23 & 2.23 & 1.284× & 4.46 & --- \emph{(rich is worse)} & \textbf{11} / 22 & 0 \\
red·B0 & 203 & 7.39 & 14.78 & 1.126× & 18.23 & yes & 0 / 0 & 0 \\
red·B1 & 203 & 2.46 & 7.39 & 1.527× & 8.37 & yes & 0 / 0 & 2 \\
red·B2 & 203 & 1.97 & 0.99 & 1.633× & 2.96 & --- \emph{(rich is worse)} & \textbf{41} / 5 & 1 \\
WA·B0 & 104 & 19.23 & 22.12 & 1.054× & 29.81 & yes & 0 / 0 & 0 \\
WA·B1 & 104 & 9.62 & 13.46 & 1.778× & 16.35 & yes & 0 / 3 & 1 \\
\bottomrule
\end{tabular}
\end{table}

% Table source-order #21: Cascade signal against a random-escalation control. Is the cascade's s
\begin{table}[tp]
\centering
\caption{Cascade signal against a random-escalation control. Is the cascade's signal doing anything? Table~\ref{tab:t20} answers a deployment question: no operating point beats always-rich on success rate: which on its own cannot distinguish a signal that carries nothing from one that carries something insufficient. This is the comparator that separates them: the same escalation budget spent at random. \textbf{Caution:} Read the median column, not the best column. The best margin is positive in 24 of 24 (cell × fraction) combinations, but each of those is a maximum over that cell's 8--10 candidate signals against a constant (the random comparator does not depend on which signal is used) so positivity there is partly arithmetic. Over all 222 (signal × cell × fraction) points, 64.4\% are positive with a median of +0.327pp: that is the unselected statement, and it is the one to quote. On cls·B2, red·B0, WA·B1 the median is negative at one or more fractions while the best is positive: there the apparent win is entirely selection. The right-hand column is how much of the gap to a per-task oracle the best signal recovers, and inherits the same caveat. The offline-splice caveat on Table~\ref{tab:t20} applies unchanged. Source: \texttt{confidence\_cascade\_with\_wa.json}.}
\label{tab:t21}
\footnotesize
\setlength{\tabcolsep}{3pt}
\begin{tabular}{@{}
  >{\raggedright\arraybackslash}p{(\linewidth - 10\tabcolsep) * \real{0.1667}}
  >{\raggedright\arraybackslash}p{(\linewidth - 10\tabcolsep) * \real{0.1667}}
  >{\raggedright\arraybackslash}p{(\linewidth - 10\tabcolsep) * \real{0.1667}}
  >{\raggedright\arraybackslash}p{(\linewidth - 10\tabcolsep) * \real{0.1667}}
  >{\raggedright\arraybackslash}p{(\linewidth - 10\tabcolsep) * \real{0.1667}}
  >{\raggedright\arraybackslash}p{(\linewidth - 10\tabcolsep) * \real{0.1667}}@{}}
\toprule
\begin{minipage}[b]{\linewidth}\raggedright
cell
\end{minipage} & \begin{minipage}[b]{\linewidth}\raggedright
signals
\end{minipage} & \begin{minipage}[b]{\linewidth}\raggedright
\textbf{best} margin (10/20/30\%)
\end{minipage} & \begin{minipage}[b]{\linewidth}\raggedright
\textbf{median over all signals} (10/20/30\%)
\end{minipage} & \begin{minipage}[b]{\linewidth}\raggedright
signals \textgreater0
\end{minipage} & \begin{minipage}[b]{\linewidth}\raggedright
oracle headroom captured (10/20/30\%)
\end{minipage} \\
\midrule
cls·B0 & 8 & +1.12 / +1.78 / +2.01 & +1.12 / +1.11 / +0.67 & 19/24 & 15\% / 25\% / 30\% \\
cls·B1 & 10 & +1.16 / +1.43 / +2.14 & +0.27 / +0.53 / +0.81 & 19/30 & 20\% / 27\% / 40\% \\
cls·B2 & 10 & +0.45 / +0.89 / +0.45 & +0.00 / +0.00 / +0.00 & 11/30 & 20\% / 40\% / 20\% \\
red·B0 & 10 & +1.74 / +1.46 / +0.74 & +0.01 / +0.48 / -0.25 & 15/30 & 23\% / 27\% / 27\% \\
red·B1 & 8 & +1.49 / +1.96 / +1.48 & +0.50 / +1.22 / +1.23 & 23/24 & 33\% / 50\% / 50\% \\
red·B2 & 9 & +0.10 / +0.20 / +0.30 & +0.10 / +0.20 / +0.30 & 20/27 & 0\% / 0\% / 0\% \\
WA·B0 & 10 & +3.57 / +3.26 / +3.95 & +1.16 / +1.82 / +1.54 & 24/30 & 36\% / 36\% / 45\% \\
WA·B1 & 9 & +2.51 / +2.11 / +2.70 & -0.37 / -0.78 / -0.18 & 12/27 & 43\% / 43\% / 57\% \\
\bottomrule
\end{tabular}
\end{table}

% Table source-order #22: Triage learnability and the visual-difficulty feature. Can the triage 
\begin{table}[tp]
\centering
\caption{Triage learnability and the visual-difficulty feature. Can the triage label be predicted, and does the benchmark's own visual-difficulty annotation help? Task-held-out 5-fold CV, L2 logistic regression, seed 42. Mean ΔAUROC = +0.0024 over 6 cells, improving 1: inside fold-split noise. \texttt{solvable} is the positive count: the label exists for every task, unlike the which-mode label. \textbf{Caution:} Six cells, not eight, and the reason is the finding. \texttt{visual\_difficulty} is a VisualWebArena task-config annotation; WebArena's 106 reddit configs carry none of \texttt{visual\_difficulty}, \texttt{reasoning\_difficulty}, \texttt{overall\_difficulty} or \texttt{image}: the field simply is not there, so the contrast cannot be computed rather than computing to zero. That absence is what makes WA the clean test elsewhere: where a router can read the benchmark's own difficulty annotation it looks learnable, and WA is the setting where it cannot. Source: \texttt{visual\_difficulty\_router.json}.}
\label{tab:t22}
\footnotesize
\setlength{\tabcolsep}{3pt}
\begin{tabular}{@{}
  >{\raggedright\arraybackslash}p{(\linewidth - 12\tabcolsep) * \real{0.1429}}
  >{\raggedright\arraybackslash}p{(\linewidth - 12\tabcolsep) * \real{0.1429}}
  >{\raggedright\arraybackslash}p{(\linewidth - 12\tabcolsep) * \real{0.1429}}
  >{\raggedright\arraybackslash}p{(\linewidth - 12\tabcolsep) * \real{0.1429}}
  >{\raggedright\arraybackslash}p{(\linewidth - 12\tabcolsep) * \real{0.1429}}
  >{\raggedright\arraybackslash}p{(\linewidth - 12\tabcolsep) * \real{0.1429}}
  >{\raggedright\arraybackslash}p{(\linewidth - 12\tabcolsep) * \real{0.1429}}@{}}
\toprule
\begin{minipage}[b]{\linewidth}\raggedright
cell
\end{minipage} & \begin{minipage}[b]{\linewidth}\raggedright
n
\end{minipage} & \begin{minipage}[b]{\linewidth}\raggedright
solvable
\end{minipage} & \begin{minipage}[b]{\linewidth}\raggedright
AUROC without
\end{minipage} & \begin{minipage}[b]{\linewidth}\raggedright
AUROC with visual\_\allowbreak{}difficulty
\end{minipage} & \begin{minipage}[b]{\linewidth}\raggedright
Δ
\end{minipage} & \begin{minipage}[b]{\linewidth}\raggedright
visual\_\allowbreak{}difficulty alone
\end{minipage} \\
\midrule
cls·B0 & 224 & 97 & 0.726 & 0.726 & -0.000 & 0.534 \\
cls·B1 & 224 & 55 & 0.732 & 0.726 & -0.006 & 0.553 \\
cls·B2 & 224 & 16 & 0.642 & 0.630 & -0.013 & 0.570 \\
red·B0 & 203 & 53 & 0.780 & 0.776 & -0.005 & 0.630 \\
red·B1 & 203 & 24 & 0.864 & 0.863 & -0.002 & 0.620 \\
red·B2 & 203 & 15 & 0.615 & 0.655 & +0.040 & 0.718 \\
\bottomrule
\end{tabular}
\end{table}

% Table source-order #23: The intuitive routing feature. The intuitive routing feature. A task s
\begin{table}[tp]
\centering
\caption{The intuitive routing feature. The intuitive routing feature. A task shipping a reference image ought to route to a mode that can see images: the table shows which mode is actually best on each side of that split. \textbf{Caution:} Reference images are delivered in every mode, so this feature does not separate what it appears to. WebArena cannot arbitrate: its 106 reddit configs carry no \texttt{image} field at all (nor any of the three difficulty annotations) so both of this table's stratifiers are undefined there. Six cells is the whole population for this question, not a coverage gap. Source: \texttt{routing\_feature\_diagnostics.json}.}
\label{tab:t23}
\footnotesize
\setlength{\tabcolsep}{3pt}
\begin{tabular}{@{}
  >{\raggedright\arraybackslash}p{(\linewidth - 12\tabcolsep) * \real{0.1429}}
  >{\raggedright\arraybackslash}p{(\linewidth - 12\tabcolsep) * \real{0.1429}}
  >{\raggedright\arraybackslash}p{(\linewidth - 12\tabcolsep) * \real{0.1429}}
  >{\raggedright\arraybackslash}p{(\linewidth - 12\tabcolsep) * \real{0.1429}}
  >{\raggedright\arraybackslash}p{(\linewidth - 12\tabcolsep) * \real{0.1429}}
  >{\raggedright\arraybackslash}p{(\linewidth - 12\tabcolsep) * \real{0.1429}}
  >{\raggedright\arraybackslash}p{(\linewidth - 12\tabcolsep) * \real{0.1429}}@{}}
\toprule
\begin{minipage}[b]{\linewidth}\raggedright
cell
\end{minipage} & \begin{minipage}[b]{\linewidth}\raggedright
tasks with ref image
\end{minipage} & \begin{minipage}[b]{\linewidth}\raggedright
without
\end{minipage} & \begin{minipage}[b]{\linewidth}\raggedright
best mode WITH image
\end{minipage} & \begin{minipage}[b]{\linewidth}\raggedright
SR
\end{minipage} & \begin{minipage}[b]{\linewidth}\raggedright
best mode WITHOUT
\end{minipage} & \begin{minipage}[b]{\linewidth}\raggedright
SR
\end{minipage} \\
\midrule
cls·B0 & 65 & 159 & SoM & 40.00 & Vision & 22.64 \\
cls·B1 & 65 & 159 & DOM+\allowbreak{}stext & 16.92 & SoM & 13.84 \\
cls·B2 & 65 & 159 & SoM & 4.62 & DOM & 1.89 \\
red·B0 & 79 & 124 & SoM & 26.58 & DOM & 7.26 \\
red·B1 & 79 & 124 & SoM & 13.92 & SoM & 3.23 \\
red·B2 & 79 & 124 & DOM & 8.86 & Vision & 3.23 \\
\bottomrule
\end{tabular}
\end{table}

% Table source-order #33: generated by scripts/analysis/export_ablation_tables.py
% Do not hand-edit the body; rerun the producer. Break-point fixes are the exception.
\begin{table}[tp]
\centering
\caption{Pooled tier router. Pooling backbones restores label supply and does not buy a router. Labels are scarce per cell, so this arm pools backbones that share a task and routes by the coarsest label available (which of two cost tiers to spend) giving classifieds 152/reddit 77/wa\_reddit 86 labelled rows against the 15--97 a single cell affords. H-pool NOT supported: the same-family × cost-tier router dominates always-cheapest in 0/6 cells and is dominated by the fixed-mode menu in every cell. The most favourable corner (agreeing backbones, coarse label, highest ceiling) does not change the negative result. Δ SR is not the story: the cost column is. Where the router does raise success it also raises cost by 2.7--55.7\%, so it is buying success at a price rather than routing to a better arm, and it strictly dominates always-cheapest in 0 of 14 arms. \textbf{Caution: non-dominance is not dominance.} The \texttt{non-dominated} column is an admissibility criterion (a policy buying success for more money passes it) and it must never be quoted as a win rate. \textbf{Caution: ‡ marks rows that are identical by construction, not by result}: WebArena has no cross-family backbone, so its all-three pool and its same-family pool are the same pool, reported twice to keep the layout parallel. \textbf{Caution:} the pooled task overlap is 50 tasks on classifieds and 20 on reddit, so every per-cell contrast is underpowered; costs are never compared across backbones (an API bill against an electricity estimate); and this probe is post-hoc and exploratory: it is not the preregistered gate and must not be cited as one. Source: \texttt{router\_pooled\_tier\_learnability.json}.}
\label{tab:t33}
\footnotesize
\setlength{\tabcolsep}{3pt}
\begin{tabular}{@{}
  >{\raggedright\arraybackslash}p{(\linewidth - 12\tabcolsep) * \real{0.1429}}
  >{\raggedright\arraybackslash}p{(\linewidth - 12\tabcolsep) * \real{0.1429}}
  >{\raggedright\arraybackslash}p{(\linewidth - 12\tabcolsep) * \real{0.1429}}
  >{\raggedright\arraybackslash}p{(\linewidth - 12\tabcolsep) * \real{0.1429}}
  >{\raggedright\arraybackslash}p{(\linewidth - 12\tabcolsep) * \real{0.1429}}
  >{\raggedright\arraybackslash}p{(\linewidth - 12\tabcolsep) * \real{0.1429}}
  >{\raggedright\arraybackslash}p{(\linewidth - 12\tabcolsep) * \real{0.1429}}@{}}
\toprule
\begin{minipage}[b]{\linewidth}\raggedright
pool
\end{minipage} & \begin{minipage}[b]{\linewidth}\raggedright
cell
\end{minipage} & \begin{minipage}[b]{\linewidth}\raggedright
router SR
\end{minipage} & \begin{minipage}[b]{\linewidth}\raggedright
cheapest SR
\end{minipage} & \begin{minipage}[b]{\linewidth}\raggedright
Δ SR
\end{minipage} & \begin{minipage}[b]{\linewidth}\raggedright
Δ cost
\end{minipage} & \begin{minipage}[b]{\linewidth}\raggedright
non-dom.
\end{minipage} \\
\midrule
same-family & \texttt{cls·B0} & 25.00\% & 25.00\% & +0.00pp & +13.9\% & 46.8\% \\
same-family & \texttt{cls·B1} & 12.05\% & 12.50\% & -0.45pp & +33.0\% & 33.6\% \\
all-three & \texttt{cls·B0} & 24.11\% & 25.00\% & -0.89pp & +14.6\% & 34.5\% \\
all-three & \texttt{cls·B1} & 11.16\% & 12.50\% & -1.34pp & +36.6\% & 21.2\% \\
all-three & \texttt{cls·B2} & 1.79\% & 2.23\% & -0.45pp & +4.8\% & 18.0\% \\
same-family & \texttt{red·B0} & 14.29\% & 7.39\% & +6.90pp & +10.2\% & 99.2\% \\
same-family & \texttt{red·B1} & 4.43\% & 2.46\% & +1.97pp & +38.2\% & 86.8\% \\
all-three & \texttt{red·B0} & 13.30\% & 7.39\% & +5.91pp & +2.7\% & 98.2\% \\
all-three & \texttt{red·B1} & 4.43\% & 2.46\% & +1.97pp & +35.8\% & 86.8\% \\
all-three & \texttt{red·B2} & 1.97\% & 1.97\% & +0.00pp & +7.4\% & 42.5\% \\
same-family & \texttt{wa\_red·B0} & 30.77\% & 26.92\% & +3.85pp & +6.3\% & 83.3\% \\
same-family & \texttt{wa\_red·B1} & 15.38\% & 9.62\% & +5.77pp & +55.7\% & 94.1\% \\
all-three ‡ & \texttt{wa\_red·B0} & 30.77\% & 26.92\% & +3.85pp & +6.3\% & 83.3\% \\
all-three ‡ & \texttt{wa\_red·B1} & 15.38\% & 9.62\% & +5.77pp & +55.7\% & 94.1\% \\
\bottomrule
\end{tabular}
\end{table}


% Paper B table #06: \caption{The nested triage policy against the always-cheapest fixed policy. Pare
\begin{table}[tp]
\centering
\caption{The nested triage policy against the always-cheapest fixed policy. Pareto dominance requires no worse on both axes; no cell achieves it. The parenthetical gives the router's deltas relative to the fixed policy.}
\label{tab:pb06}
\footnotesize
\setlength{\tabcolsep}{3pt}
\begin{tabular}{@{}
  >{\raggedright\arraybackslash}p{(\linewidth - 6\tabcolsep) * \real{0.2500}}
  >{\raggedright\arraybackslash}p{(\linewidth - 6\tabcolsep) * \real{0.2500}}
  >{\raggedright\arraybackslash}p{(\linewidth - 6\tabcolsep) * \real{0.2500}}
  >{\raggedright\arraybackslash}p{(\linewidth - 6\tabcolsep) * \real{0.2500}}@{}}
\toprule
\begin{minipage}[b]{\linewidth}\raggedright
cell
\end{minipage} & \begin{minipage}[b]{\linewidth}\raggedright
nested SR / cost
\end{minipage} & \begin{minipage}[b]{\linewidth}\raggedright
always-cheapest SR / cost
\end{minipage} & \begin{minipage}[b]{\linewidth}\raggedright
Pareto-dominates?
\end{minipage} \\
\midrule
classifieds · B0 & 26.79\% / 0.07197 & 25.00\% / 0.06481 & no (+1.79pp SR, +11.0\% cost) \\
reddit · B0 & 12.81\% / 0.09836 & 7.39\% / 0.09807 & no (+5.42pp SR, +0.3\% cost) \\
classifieds · B1 & 14.29\% / 0.05757 & 12.50\% / 0.04316 & no (+1.79pp SR, +33.4\% cost) \\
reddit · B1 & 5.91\% / 0.06970 & 2.46\% / 0.05240 & no (+3.45pp SR, +33.0\% cost) \\
classifieds · B2 & 1.34\% / 0.07247 & 2.23\% / 0.07065 & no (−0.89pp SR, +2.6\% cost) \\
reddit · B2 & 3.94\% / 0.06964 & 1.97\% / 0.06833 & no (+1.97pp SR, +1.9\% cost) \\
\bottomrule
\end{tabular}
\end{table}

% Paper B table #07: \caption{The five supervision targets against the two requirements a trainable r
\begin{table}[tp]
\centering
\caption{The five supervision targets against the two requirements a trainable router needs. Each closed route is closed for a different reason, which is what makes the negative result closed rather than provisional. Supply and identifiability are judged on that target's own denominator, which differs between the last two rows.}
\label{tab:pb07}
\footnotesize
\setlength{\tabcolsep}{3pt}
\begin{tabular}{@{}
  >{\raggedright\arraybackslash}p{(\linewidth - 6\tabcolsep) * \real{0.2500}}
  >{\raggedright\arraybackslash}p{(\linewidth - 6\tabcolsep) * \real{0.2500}}
  >{\raggedright\arraybackslash}p{(\linewidth - 6\tabcolsep) * \real{0.2500}}
  >{\raggedright\arraybackslash}p{(\linewidth - 6\tabcolsep) * \real{0.2500}}@{}}
\toprule
\begin{minipage}[b]{\linewidth}\raggedright
route
\end{minipage} & \begin{minipage}[b]{\linewidth}\raggedright
supply
\end{minipage} & \begin{minipage}[b]{\linewidth}\raggedright
identifiability
\end{minipage} & \begin{minipage}[b]{\linewidth}\raggedright
outcome
\end{minipage} \\
\midrule
which-mode, per cell & fails (15--97 labels, 4/6 cells untrainable) & fine & closed by \S\ref{sec:lowerbound} \\
continuous target & would fix & fine & closed: benchmark score is binary \\
pooled which-mode & fixed (260) & fails (56--57\% conflict; in-sample modal agreement 79--84\%) & wrong estimand \\
triage (binary) & fine (203--224, every task) & fine & learnable, and beaten by a fixed policy (\S\ref{sec:lowerbound}) \\
screenshot tier & fine, but only on solved tasks & fine (68--88\% cross-backbone agreement, over solved tasks) & agreement measured; no tier classifier trained \\
\bottomrule
\end{tabular}
\end{table}


\section{Failure attribution and mechanisms}\label{app:failures}

Table~\ref{tab:t43} locates \emph{when} two modes part on the tasks whose outcome they disagree about: at the first or second step, essentially never late, which is not the shape accumulated drift has. It is markedly weaker on WebArena, so it is reported per site rather than pooled. Table~\ref{tab:t41} is the paired cut: on tasks only one channel solved, how did the other fail? It counts rule hits, so Table~\ref{tab:t26} asks the same question without the rule vocabulary, using six probes read straight off the step records, and finds the losing channel failing \emph{more blandly} there than elsewhere, which is what closes the objection that the residual is an artefact of a VWA-shaped ruleset. Table~\ref{tab:t30} decomposes the DOM$\rightarrow$SoM-image transition into a text axis and a prompt axis, and Table~\ref{tab:t31} asks whether that text axis moves per-step decision quality more than it moves macro action frequencies. Table~\ref{tab:t32} counts hallucinated element references, which are inapplicable to Vision by construction, and the table marks them. \textbf{Caution:} a per-rule frequency is a distribution of symptoms, not of causes. The two largest rows in most cells are risk markers that causal verification did not confirm as death causes.


Each block of Table~\ref{tab:t41} states the largest enrichment among the rules it does \emph{not} list. Read each block against its own baseline column, never across blocks: the text channel is four arms and the image channel two, so the two blocks' task counts are not comparable to each other.

% Table source-order #26: Vocabulary-free probes on the text-wins residual. Is the text-wins res
\begin{table}[tp]
\centering
\caption{Vocabulary-free probes on the text-wins residual. Is the text-wins residual real, or an artefact of the rule vocabulary? Tables 24--25 count rule hits, and the ruleset was discovered on VisualWebArena: so an absent signature there could be a property of the vocabulary rather than of the world. Each probe here is computed from raw step fields and never from a rule hit, over 142 disagreement episodes against 2304 baseline failures. The largest enrichment is 1.15× and 5 of 6 sit \emph{below} 1: on the tasks the text channel uniquely solves, the image channel fails more blandly than it fails elsewhere: it did not arrive, rather than breaking somewhere nameable. \textbf{Caution:} Six candidates chosen by us, so this cannot show that no mechanism exists; what it closes is the specific objection that the residual is an artefact of a VWA-shaped vocabulary. Source: \texttt{conditional\_failure\_attribution.json}}
\label{tab:t26}
\footnotesize
\setlength{\tabcolsep}{3pt}
\begin{tabular}{@{}
  >{\raggedright\arraybackslash}p{(\linewidth - 6\tabcolsep) * \real{0.2500}}
  >{\raggedright\arraybackslash}p{(\linewidth - 6\tabcolsep) * \real{0.2500}}
  >{\raggedright\arraybackslash}p{(\linewidth - 6\tabcolsep) * \real{0.2500}}
  >{\raggedright\arraybackslash}p{(\linewidth - 6\tabcolsep) * \real{0.2500}}@{}}
\toprule
\begin{minipage}[b]{\linewidth}\raggedright
candidate mechanism
\end{minipage} & \begin{minipage}[b]{\linewidth}\raggedright
on the disagreement set
\end{minipage} & \begin{minipage}[b]{\linewidth}\raggedright
that channel's baseline
\end{minipage} & \begin{minipage}[b]{\linewidth}\raggedright
enrichment
\end{minipage} \\
\midrule
finished in five steps or fewer & 0.162 & 0.141 & \textbf{1.15×} \\
any parse failure & 0.092 & 0.109 & \textbf{0.84×} \\
never searched & 0.387 & 0.468 & \textbf{0.83×} \\
ran out of budget (\textgreater=30 steps) & 0.535 & 0.658 & \textbf{0.81×} \\
page unchanged on over half the steps & 0.331 & 0.431 & \textbf{0.77×} \\
action failure on over half the steps & 0.310 & 0.416 & \textbf{0.74×} \\
\bottomrule
\end{tabular}
\end{table}

% Table source-order #41: Failure modes are asymmetric across channels.
\begin{table}[tbp]
\centering
\caption{Failure modes are asymmetric across channels. The two channels do not fail the same way. On tasks only one channel solved, how did the other fail? Pooled over 8 cells at ruleset v11. Enrichment = hit rate on the disagreement set over that channel's hit rate across all its failures, so about 1x means it failed there the way it fails everywhere. Rules with fewer than 8 pooled conditional hits are omitted, and the largest omitted enrichment is stated per block rather than left to the reader. Block A has named death causes (top 2.31x, 404 losing-channel episodes against 5388 of that channel's failures overall). Block B does not: its top row rests on 8 hits, exactly the reporting floor, and all 8 of them fall in the 2 WebArena cells, none in the other 6; every other rule in that block sits at or below the everywhere-baseline. On the tasks the text channel uniquely solves, the image channel did not break somewhere nameable -- it did not arrive. \textbf{Caution:} TEXT is four arms and IMAGE is two, so the two blocks' task counts are not comparable to each other; read each block against its own baseline column, never across blocks. \textbf{Caution:} a per-rule frequency is a distribution of symptoms, not of causes -- the largest rows in most cells are risk markers, not death causes, and only rules whose docstrings record a causal check are verified as such. Source: \texttt{conditional\_failure\_attribution.json}.}
\label{tab:t41}
\small
\setlength{\tabcolsep}{4pt}
\begin{tabular}{@{}
  >{\raggedright\arraybackslash}p{(\linewidth - 12\tabcolsep) * \real{0.1429}}
  >{\raggedright\arraybackslash}p{(\linewidth - 12\tabcolsep) * \real{0.1429}}
  >{\raggedright\arraybackslash}p{(\linewidth - 12\tabcolsep) * \real{0.1429}}
  >{\raggedright\arraybackslash}p{(\linewidth - 12\tabcolsep) * \real{0.1429}}
  >{\raggedright\arraybackslash}p{(\linewidth - 12\tabcolsep) * \real{0.1429}}
  >{\raggedright\arraybackslash}p{(\linewidth - 12\tabcolsep) * \real{0.1429}}
  >{\raggedright\arraybackslash}p{(\linewidth - 12\tabcolsep) * \real{0.1429}}@{}}
\toprule
\begin{minipage}[b]{\linewidth}\raggedright
\end{minipage} & \begin{minipage}[b]{\linewidth}\raggedright
rule
\end{minipage} & \begin{minipage}[b]{\linewidth}\raggedright
how it failed
\end{minipage} & \begin{minipage}[b]{\linewidth}\raggedright
on disagreement
\end{minipage} & \begin{minipage}[b]{\linewidth}\raggedright
baseline
\end{minipage} & \begin{minipage}[b]{\linewidth}\raggedright
enrichment
\end{minipage} & \begin{minipage}[b]{\linewidth}\raggedright
hits
\end{minipage} \\
\midrule
\textbf{A. image wins} & & \emph{how the text channel failed} & (101 tasks) & & & \\
& \texttt{P27} & gives up when not found & 3.2\% & 1.4\% & \textbf{2.31x} & 13 \\
& \texttt{P17} & click-back oscillation & 15.1\% & 6.7\% & \textbf{2.25x} & 61 \\
& \texttt{P16} & visual task, DOM cannot see & 6.2\% & 2.8\% & \textbf{2.24x} & 25 \\
& \texttt{P43} & page visual, no screenshot & 48.5\% & 29.3\% & \textbf{1.65x} & 196 \\
& & \emph{15 further rules} & & & all \textless= 1.55x & \\
\textbf{B. text wins} & & \emph{how the image channel failed} & (109 tasks) & & & \\
& \texttt{P49} & submit-page anchor misclick & 3.7\% & 1.0\% & \textbf{3.61x} & 8 \\
& \texttt{P17} & click-back oscillation & 4.6\% & 3.9\% & \textbf{1.17x} & 10 \\
& \texttt{P12} & never paginates & 13.8\% & 14.8\% & \textbf{0.93x} & 30 \\
& \texttt{P31} & budget exhausted, unfinished & 49.5\% & 54.2\% & \textbf{0.91x} & 108 \\
& & \emph{6 further rules} & & & all \textless= 0.87x & \\
\bottomrule
\end{tabular}
\end{table}


% Table source-order #43: Where two representations part company (E2 trajectory divergence).
\begin{table}[tbp]
\centering
\caption{Where two representations part company. Where two representations part company. Restricted to the tasks whose outcome the two modes disagree on (exactly one of them succeeds) so these are the trajectories that produced the difference, not all trajectories. \texttt{first\ divergent\ step} is the first step at which the two URL signatures differ. Drift would look different: two runs that separate through accumulated perturbation share a prefix and part late, whereas on VisualWebArena 72--100\% of these tasks have already diverged by step 3 and at most 3.0\% diverge at step 10 or later. \textbf{Caution: this is markedly weaker on WebArena} (50--89\% early, up to 22.2\% late), which is why the table is per site and not pooled. \textbf{Caution:} each row rests on 13--46 tasks, and element ids are deliberately not used as the anchor: they are neither step-invariant nor mode-invariant, so URL signatures are the only stable comparison. Source: \texttt{mechanism\_per\_task.json} E2.}
\label{tab:t43}
\small
\setlength{\tabcolsep}{4pt}
\begin{tabular}{@{}
  >{\raggedright\arraybackslash}p{(\linewidth - 10\tabcolsep) * \real{0.1667}}
  >{\raggedright\arraybackslash}p{(\linewidth - 10\tabcolsep) * \real{0.1667}}
  >{\raggedright\arraybackslash}p{(\linewidth - 10\tabcolsep) * \real{0.1667}}
  >{\raggedright\arraybackslash}p{(\linewidth - 10\tabcolsep) * \real{0.1667}}
  >{\raggedright\arraybackslash}p{(\linewidth - 10\tabcolsep) * \real{0.1667}}
  >{\raggedright\arraybackslash}p{(\linewidth - 10\tabcolsep) * \real{0.1667}}@{}}
\toprule
\begin{minipage}[b]{\linewidth}\raggedright
site
\end{minipage} & \begin{minipage}[b]{\linewidth}\raggedright
contrast
\end{minipage} & \begin{minipage}[b]{\linewidth}\raggedright
tasks
\end{minipage} & \begin{minipage}[b]{\linewidth}\raggedright
median first divergent step
\end{minipage} & \begin{minipage}[b]{\linewidth}\raggedright
diverged by step 3
\end{minipage} & \begin{minipage}[b]{\linewidth}\raggedright
diverged at step 10+
\end{minipage} \\
\midrule
classifieds & DOM vs SoM-image & 22 & 1 & 95.5\% & 0.0\% \\
classifieds & DOM vs DOM+\allowbreak{}sprompt & 33 & 1 & 84.8\% & 3.0\% \\
classifieds & DOM vs DOM+\allowbreak{}stext & 22 & 1 & 95.5\% & 0.0\% \\
classifieds & SoM-image vs SoM & 46 & 1 & 93.5\% & 0.0\% \\
classifieds & DOM+\allowbreak{}sprompt vs SoM-image & 23 & 1 & 91.3\% & 0.0\% \\
classifieds & DOM+\allowbreak{}stext vs SoM-image & 18 & 2 & 72.2\% & 0.0\% \\
reddit & DOM vs SoM-image & 19 & 0 & 100.0\% & 0.0\% \\
reddit & DOM vs DOM+\allowbreak{}sprompt & 18 & 1 & 88.9\% & 0.0\% \\
reddit & DOM vs DOM+\allowbreak{}stext & 16 & 0 & 93.8\% & 0.0\% \\
reddit & SoM-image vs SoM & 20 & 0 & 94.7\% & 0.0\% \\
reddit & DOM+\allowbreak{}sprompt vs SoM-image & 17 & 0 & 100.0\% & 0.0\% \\
reddit & DOM+\allowbreak{}stext vs SoM-image & 19 & 0 & 100.0\% & 0.0\% \\
wa\_reddit & DOM vs SoM-image & 16 & 1 & 75.0\% & 6.2\% \\
wa\_reddit & DOM vs DOM+\allowbreak{}sprompt & 19 & 1 & 89.5\% & 5.3\% \\
wa\_reddit & DOM vs DOM+\allowbreak{}stext & 21 & 1 & 66.7\% & 9.5\% \\
wa\_reddit & SoM-image vs SoM & 19 & 1 & 73.7\% & 0.0\% \\
wa\_reddit & DOM+\allowbreak{}sprompt vs SoM-image & 13 & 0.5 & 83.3\% & 8.3\% \\
wa\_reddit & DOM+\allowbreak{}stext vs SoM-image & 19 & 4 & 50.0\% & 22.2\% \\
\bottomrule
\end{tabular}
\end{table}


% Table source-order #30: 2x2 axis decomposition. 2×2 axis decomposition. Effect sizes are Cohen
\begin{table}[tp]
\centering
\caption{2x2 axis decomposition. 2×2 axis decomposition. Effect sizes are Cohen's h (binary metrics) or d\_z (paired continuous), signed right-minus-left. The compound DOM→SoM-image transition decomposes into a text-payload axis and a prompt-style axis; the image axis is SoM-image→SoM. \textbf{Caution:} On mean differences the two decomposition routes agreeing is an algebraic identity, so a zero residual is arithmetic and not evidence about an interaction. \texttt{B2\ ×\ wa\_reddit} is absent because B2 never ran WebArena. Source: \texttt{axis\_effect\_size\_with\_wa.json}.}
\label{tab:t30}
\scriptsize
\setlength{\tabcolsep}{3pt}
\begin{tabular}{@{}llllll@{}}
\toprule
cell & metric & text axis & prompt axis & image axis & DOM→SoM-image \\
\midrule
cls·B0 & search\_loop & -0.011 & -0.045 & -0.225 & -0.056 \\
cls·B0 & type\_frac & -0.169 & +0.049 & +0.030 & -0.140 \\
cls·B0 & scroll\_frac & +0.072 & +0.063 & -0.253 & +0.134 \\
cls·B0 & selfcorr\_count & +0.088 & -0.115 & -0.091 & -0.020 \\
cls·B0 & click\_frac & -0.001 & -0.059 & +0.035 & -0.056 \\
cls·B0 & finish\_rate & -0.070 & -0.039 & +0.119 & -0.109 \\
cls·B0 & n\_steps & +0.023 & +0.039 & -0.246 & +0.066 \\
cls·B0 & action\_repeat\_frac & -0.051 & +0.086 & -0.117 & +0.030 \\
cls·B1 & search\_loop & +0.066 & -0.160 & -0.203 & -0.094 \\
cls·B1 & type\_frac & +0.088 & -0.546 & -0.005 & -0.452 \\
cls·B1 & scroll\_frac & -0.097 & +0.025 & -0.062 & -0.065 \\
cls·B1 & selfcorr\_count & -0.015 & +0.007 & +0.149 & -0.007 \\
cls·B1 & click\_frac & +0.139 & +0.526 & -0.169 & +0.614 \\
cls·B1 & finish\_rate & -0.055 & +0.082 & +0.170 & +0.027 \\
cls·B1 & n\_steps & +0.101 & -0.108 & -0.253 & -0.011 \\
cls·B1 & action\_repeat\_frac & +0.053 & +0.002 & -0.174 & +0.048 \\
cls·B2 & search\_loop & -0.037 & +0.113 & -0.141 & +0.076 \\
cls·B2 & type\_frac & +0.116 & -0.410 & +0.297 & -0.290 \\
cls·B2 & scroll\_frac & +0.154 & -0.112 & +0.003 & +0.026 \\
cls·B2 & selfcorr\_count & +0.100 & -0.111 & -0.105 & -0.024 \\
cls·B2 & click\_frac & -0.476 & +0.485 & -0.097 & +0.011 \\
cls·B2 & finish\_rate & -0.217 & -0.038 & +0.333 & -0.256 \\
cls·B2 & n\_steps & -0.052 & +0.160 & -0.363 & +0.115 \\
cls·B2 & action\_repeat\_frac & -0.093 & +0.159 & -0.067 & +0.064 \\
red·B0 & search\_loop & -0.170 & -0.104 & +0.052 & -0.275 \\
red·B0 & type\_frac & +0.029 & -0.098 & -0.084 & -0.067 \\
red·B0 & scroll\_frac & -0.206 & +0.041 & -0.148 & -0.182 \\
red·B0 & selfcorr\_count & +0.279 & -0.065 & -0.258 & +0.223 \\
red·B0 & click\_frac & -0.057 & +0.053 & +0.081 & -0.009 \\
red·B0 & finish\_rate & -0.299 & +0.102 & +0.160 & -0.200 \\
red·B0 & n\_steps & +0.276 & -0.032 & -0.252 & +0.245 \\
red·B0 & action\_repeat\_frac & +0.226 & -0.056 & -0.033 & +0.202 \\
red·B1 & search\_loop & +0.020 & -0.100 & -0.297 & -0.080 \\
red·B1 & type\_frac & +0.032 & -0.329 & -0.159 & -0.301 \\
red·B1 & scroll\_frac & +0.008 & -0.071 & +0.157 & -0.074 \\
red·B1 & selfcorr\_count & -0.053 & -0.075 & -0.062 & -0.120 \\
red·B1 & click\_frac & +0.019 & +0.326 & -0.016 & +0.370 \\
red·B1 & finish\_rate & -0.224 & +0.012 & +0.223 & -0.212 \\
red·B1 & n\_steps & +0.202 & -0.044 & -0.231 & +0.165 \\
red·B1 & action\_repeat\_frac & +0.133 & +0.108 & +0.020 & +0.249 \\
red·B2 & search\_loop & +0.092 & -0.115 & -0.179 & -0.024 \\
red·B2 & type\_frac & -0.045 & -0.023 & +0.045 & -0.070 \\
red·B2 & scroll\_frac & -0.037 & +0.093 & -0.040 & +0.054 \\
red·B2 & selfcorr\_count & +0.052 & -0.026 & +0.015 & +0.019 \\
red·B2 & click\_frac & -0.193 & +0.143 & -0.017 & -0.047 \\
red·B2 & finish\_rate & -0.059 & -0.044 & +0.382 & -0.103 \\
red·B2 & n\_steps & -0.126 & +0.057 & -0.137 & -0.054 \\
red·B2 & action\_repeat\_frac & -0.157 & -0.010 & +0.121 & -0.173 \\
WA·B0 & search\_loop & -0.042 & -0.155 & +0.238 & -0.197 \\
WA·B0 & type\_frac & +0.476 & -0.197 & -0.031 & +0.336 \\
WA·B0 & scroll\_frac & -0.247 & +0.155 & -0.191 & -0.135 \\
WA·B0 & selfcorr\_count & +0.106 & -0.053 & +0.088 & +0.078 \\
WA·B0 & click\_frac & -0.255 & +0.111 & -0.139 & -0.170 \\
WA·B0 & finish\_rate & -0.339 & +0.176 & +0.000 & -0.163 \\
WA·B0 & n\_steps & +0.260 & -0.084 & -0.128 & +0.200 \\
WA·B0 & action\_repeat\_frac & +0.157 & +0.012 & -0.217 & +0.170 \\
WA·B1 & search\_loop & +0.081 & -0.200 & -0.270 & -0.119 \\
WA·B1 & type\_frac & -0.016 & -0.271 & -0.109 & -0.275 \\
WA·B1 & scroll\_frac & +0.072 & -0.133 & +0.155 & -0.064 \\
WA·B1 & selfcorr\_count & +0.148 & -0.136 & +0.108 & -0.006 \\
WA·B1 & click\_frac & +0.153 & +0.244 & -0.182 & +0.353 \\
WA·B1 & finish\_rate & +0.121 & -0.204 & -0.021 & -0.083 \\
WA·B1 & n\_steps & +0.071 & +0.058 & +0.010 & +0.121 \\
WA·B1 & action\_repeat\_frac & -0.034 & +0.218 & +0.029 & +0.168 \\
\bottomrule
\end{tabular}
\end{table}

% Table source-order #31: Decision quality versus macro frequency. Does the text axis change per
\begin{table}[tp]
\centering
\caption{Decision quality versus macro frequency. Does the text axis change per-step decisions more than it changes macro action frequencies? Ratio \textgreater1 means yes. Verdict: generalizes (site\_ok = \{`reddit': True, `classifieds': True, `wa\_reddit': True\}). \textbf{Caution:} \texttt{\_site\_ok} passes a site if any backbone clears 1.0, which is a loose bar: \texttt{WA·B0} is 0.97 and the WA site passes on B1's 2.98. Until 2026-08-03 the verdict function named only the two VWA sites literally and did not consult WA at all. Source: \texttt{axis1\_microbehavior\_with\_wa.json}.}
\label{tab:t31}
\footnotesize
\setlength{\tabcolsep}{3pt}
\begin{tabular}{@{}
  >{\raggedright\arraybackslash}p{(\linewidth - 8\tabcolsep) * \real{0.2000}}
  >{\raggedright\arraybackslash}p{(\linewidth - 8\tabcolsep) * \real{0.2000}}
  >{\raggedright\arraybackslash}p{(\linewidth - 8\tabcolsep) * \real{0.2000}}
  >{\raggedright\arraybackslash}p{(\linewidth - 8\tabcolsep) * \real{0.2000}}
  >{\raggedright\arraybackslash}p{(\linewidth - 8\tabcolsep) * \real{0.2000}}@{}}
\toprule
\begin{minipage}[b]{\linewidth}\raggedright
cell
\end{minipage} & \begin{minipage}[b]{\linewidth}\raggedright
decision effect (mean abs)
\end{minipage} & \begin{minipage}[b]{\linewidth}\raggedright
macro effect (mean abs)
\end{minipage} & \begin{minipage}[b]{\linewidth}\raggedright
ratio
\end{minipage} & \begin{minipage}[b]{\linewidth}\raggedright
\textgreater1?
\end{minipage} \\
\midrule
cls·B0 & 0.1530 & 0.0606 & \textbf{2.52} & yes \\
cls·B1 & 0.1024 & 0.0766 & \textbf{1.34} & yes \\
cls·B2 & 0.3159 & 0.1556 & \textbf{2.03} & yes \\
red·B0 & 0.2761 & 0.1926 & \textbf{1.43} & yes \\
red·B1 & 0.2450 & 0.0863 & \textbf{2.84} & yes \\
red·B2 & 0.3877 & 0.0952 & \textbf{4.07} & yes \\
WA·B0 & 0.2284 & 0.2351 & \textbf{0.97} & \textbf{no} \\
WA·B1 & 0.2588 & 0.0869 & \textbf{2.98} & yes \\
\bottomrule
\end{tabular}
\end{table}

% Table source-order #32: Hallucinated element references. Hallucinated element references, rule
\begin{table}[tp]
\centering
\caption{Hallucinated element references. Hallucinated element references, ruleset \texttt{11-intent-text-fallback} over 36 conditions. An action naming an element id that is not in the observation. \textbf{Caution:} \texttt{vision} carries no element-id list at all, so this rule is structurally inapplicable there rather than measuring zero: the same gate-versus-measurement confusion flagged for P2/P4. Source: \texttt{cross\_mode\_failure\_signatures.json}.}
\label{tab:t32}
\footnotesize
\setlength{\tabcolsep}{3pt}
\begin{tabular}{@{}
  >{\raggedright\arraybackslash}p{(\linewidth - 10\tabcolsep) * \real{0.1667}}
  >{\raggedright\arraybackslash}p{(\linewidth - 10\tabcolsep) * \real{0.1667}}
  >{\raggedright\arraybackslash}p{(\linewidth - 10\tabcolsep) * \real{0.1667}}
  >{\raggedright\arraybackslash}p{(\linewidth - 10\tabcolsep) * \real{0.1667}}
  >{\raggedright\arraybackslash}p{(\linewidth - 10\tabcolsep) * \real{0.1667}}
  >{\raggedright\arraybackslash}p{(\linewidth - 10\tabcolsep) * \real{0.1667}}@{}}
\toprule
\begin{minipage}[b]{\linewidth}\raggedright
cell
\end{minipage} & \begin{minipage}[b]{\linewidth}\raggedright
mode
\end{minipage} & \begin{minipage}[b]{\linewidth}\raggedright
episodes
\end{minipage} & \begin{minipage}[b]{\linewidth}\raggedright
failed
\end{minipage} & \begin{minipage}[b]{\linewidth}\raggedright
with hallucinated ref
\end{minipage} & \begin{minipage}[b]{\linewidth}\raggedright
rate of failed
\end{minipage} \\
\midrule
classifieds·B0 & DOM & 224 & 185 & 15 & 8.1\% \\
classifieds·B0 & DOM+\allowbreak{}sprompt & 224 & 180 & 11 & 6.1\% \\
classifieds·B0 & SoM-image & 224 & 189 & 2 & 1.1\% \\
classifieds·B0 & DOM+\allowbreak{}stext & 224 & 189 & 9 & 4.8\% \\
classifieds·B0 & SoM & 187 & 129 & 2 & 1.6\% \\
classifieds·B0 & Vision & 224 & 168 & 0 & 0.0\% \\
reddit·B0 & DOM & 203 & 174 & 8 & 4.6\% \\
reddit·B0 & DOM+\allowbreak{}sprompt & 203 & 178 & 8 & 4.5\% \\
reddit·B0 & SoM-image & 203 & 181 & 1 & 0.6\% \\
reddit·B0 & DOM+\allowbreak{}stext & 203 & 176 & 6 & 3.4\% \\
reddit·B0 & SoM & 203 & 173 & 2 & 1.2\% \\
reddit·B0 & Vision & 203 & 188 & 0 & 0.0\% \\
classifieds·B1 & DOM & 224 & 210 & 20 & 9.5\% \\
classifieds·B1 & DOM+\allowbreak{}sprompt & 224 & 209 & 33 & 15.8\% \\
classifieds·B1 & SoM-image & 224 & 209 & 7 & 3.3\% \\
classifieds·B1 & DOM+\allowbreak{}stext & 224 & 207 & 45 & 21.7\% \\
classifieds·B1 & SoM & 224 & 192 & 0 & 0.0\% \\
classifieds·B1 & Vision & 224 & 196 & 0 & 0.0\% \\
reddit·B1 & DOM & 203 & 191 & 25 & 13.1\% \\
reddit·B1 & DOM+\allowbreak{}sprompt & 203 & 192 & 33 & 17.2\% \\
reddit·B1 & SoM-image & 203 & 191 & 3 & 1.6\% \\
reddit·B1 & DOM+\allowbreak{}stext & 203 & 191 & 24 & 12.6\% \\
reddit·B1 & SoM & 203 & 188 & 5 & 2.7\% \\
reddit·B1 & Vision & 203 & 198 & 0 & 0.0\% \\
classifieds·B2 & DOM & 224 & 221 & 82 & 37.1\% \\
classifieds·B2 & DOM+\allowbreak{}sprompt & 224 & 220 & 114 & 51.8\% \\
classifieds·B2 & SoM-image & 224 & 222 & 43 & 19.4\% \\
classifieds·B2 & DOM+\allowbreak{}stext & 224 & 223 & 75 & 33.6\% \\
classifieds·B2 & SoM & 224 & 219 & 39 & 17.8\% \\
classifieds·B2 & Vision & 224 & 219 & 0 & 0.0\% \\
reddit·B2 & DOM & 203 & 195 & 127 & 65.1\% \\
reddit·B2 & DOM+\allowbreak{}sprompt & 203 & 203 & 155 & 76.4\% \\
reddit·B2 & SoM-image & 203 & 202 & 62 & 30.7\% \\
reddit·B2 & DOM+\allowbreak{}stext & 203 & 199 & 51 & 25.6\% \\
reddit·B2 & SoM & 203 & 201 & 59 & 29.4\% \\
reddit·B2 & Vision & 203 & 199 & 0 & 0.0\% \\
\bottomrule
\end{tabular}
\end{table}


\section{Audits and provenance}\label{app:audits}

The dispatch-path audit (Table~\ref{tab:t13}), the off-site navigation and container-latency audit (Table~\ref{tab:t29}), the page-change detector correction (Table~\ref{tab:t34}), the evaluator granularity check (Table~\ref{tab:t35}), and the full earned-versus-leaked audit including the first WebArena pass (Table~\ref{tab:t40}).

% Table source-order #13: What actually delivered the click. How each action reached the browser
\begin{table}[tp]
\centering
\caption{What actually delivered the click. How each action reached the browser. \texttt{Vision} is on the coordinate path by construction (it emits no element ids) so its action success is capped by this harness's coordinate implementation (39\%) rather than by the 89\% the element-id path achieves. That is not a confound to remove (it is what screenshot-only \emph{is}), but the Vision arm measures our grounding code as much as the representation. Separately the element-id fallback share rises with backbone weakness: B0 12\% · B1 35\% · B2 37\% on the text arms: \emph{how often} a run falls back is a model property, the fallback's own 16\% success is ours. No success rate elsewhere is adjusted by this; it bounds external validity. Source: \texttt{dispatch\_path\_audit.json}}
\label{tab:t13}
\footnotesize
\setlength{\tabcolsep}{3pt}
\begin{tabular}{@{}lll@{}}
\toprule
delivery path & actions & action success \\
\midrule
id locator & 8,857 & \textbf{88.9\%} \\
other & 1,530 & \textbf{65.9\%} \\
coord & 2,138 & \textbf{38.6\%} \\
id framework & 4,564 & \textbf{16.1\%} \\
\bottomrule
\end{tabular}
\end{table}

% Table source-order #29: Off-site navigation and container latency. Off-site navigation. Postmi
\begin{table}[tp]
\centering
\caption{Off-site navigation and container latency. Off-site navigation. Postmill is a link aggregator, so an agent opening a trending thread can walk onto the live public internet; classifieds is self-contained. \textbf{Caution:} Off-site steps are faster, not slower (commercial CDNs beat a Postmill container sharing a host with the agent) so the distortion runs opposite to the intuition. The larger asymmetry is in the last two columns of the on-site medians: reddit's container costs \textasciitilde1.69× what classifieds' does before any agent behaviour enters, which is why no between-site latency number is quotable bare. Source: \texttt{offsite\_navigation\_audit.json}.}
\label{tab:t29}
\footnotesize
\setlength{\tabcolsep}{3pt}
\begin{tabular}{@{}
  >{\raggedright\arraybackslash}p{(\linewidth - 10\tabcolsep) * \real{0.1667}}
  >{\raggedright\arraybackslash}p{(\linewidth - 10\tabcolsep) * \real{0.1667}}
  >{\raggedright\arraybackslash}p{(\linewidth - 10\tabcolsep) * \real{0.1667}}
  >{\raggedright\arraybackslash}p{(\linewidth - 10\tabcolsep) * \real{0.1667}}
  >{\raggedright\arraybackslash}p{(\linewidth - 10\tabcolsep) * \real{0.1667}}
  >{\raggedright\arraybackslash}p{(\linewidth - 10\tabcolsep) * \real{0.1667}}@{}}
\toprule
\begin{minipage}[b]{\linewidth}\raggedright
cell
\end{minipage} & \begin{minipage}[b]{\linewidth}\raggedright
off-site steps
\end{minipage} & \begin{minipage}[b]{\linewidth}\raggedright
off-site episodes
\end{minipage} & \begin{minipage}[b]{\linewidth}\raggedright
median env\_\allowbreak{}step on-site
\end{minipage} & \begin{minipage}[b]{\linewidth}\raggedright
off-site
\end{minipage} & \begin{minipage}[b]{\linewidth}\raggedright
ratio
\end{minipage} \\
\midrule
B0·VWA-cla & 0/20646 (0.00\%) & 0/1344 (0.0\%) & 4,493 ms & --- & --- \\
B0·VWA-red & 478/26425 (1.81\%) & 36/1230 (2.9\%) & 11,349 ms & 10,007 ms & 0.88× \\
B1·VWA-cla & 0/27927 (0.00\%) & 0/1344 (0.0\%) & 5,750 ms & --- & --- \\
B1·VWA-red & 501/29309 (1.71\%) & 54/1230 (4.4\%) & 7,848 ms & 6,349 ms & 0.81× \\
B2·VWA-cla & 58/36529 (0.16\%) & 4/1344 (0.3\%) & 4,656 ms & 14,884 ms & 3.20× \\
B2·VWA-red & 719/33749 (2.13\%) & 49/1230 (4.0\%) & 9,324 ms & 4,850 ms & 0.52× \\
B1·WA-red & 278/14695 (1.89\%) & 40/624 (6.4\%) & 6,809 ms & 6,105 ms & 0.90× \\
B0·WA-red & 123/11703 (1.05\%) & 21/624 (3.4\%) & 6,613 ms & 11,405 ms & 1.72× \\
\bottomrule
\end{tabular}
\end{table}

% Table source-order #34: page\_changed false positives. \texttt{page\_changed} false positives.
\begin{table}[tp]
\centering
\caption{page\_changed false positives. \texttt{page\_changed} false positives. A step can register a page change that is purely cosmetic; correcting for it raises every mode's no-change rate. The Micro conclusion is unaffected (Vision remains highest in every cell either way) but a router firing on a 2-step no-change streak would trigger 5321 → 5910 times (+11.1\%). Source: \texttt{page\_change\_corrected.json}.}
\label{tab:t34}
\footnotesize
\setlength{\tabcolsep}{3pt}
\begin{tabular}{@{}
  >{\raggedright\arraybackslash}p{(\linewidth - 8\tabcolsep) * \real{0.2000}}
  >{\raggedright\arraybackslash}p{(\linewidth - 8\tabcolsep) * \real{0.2000}}
  >{\raggedright\arraybackslash}p{(\linewidth - 8\tabcolsep) * \real{0.2000}}
  >{\raggedright\arraybackslash}p{(\linewidth - 8\tabcolsep) * \real{0.2000}}
  >{\raggedright\arraybackslash}p{(\linewidth - 8\tabcolsep) * \real{0.2000}}@{}}
\toprule
\begin{minipage}[b]{\linewidth}\raggedright
mode
\end{minipage} & \begin{minipage}[b]{\linewidth}\raggedright
no-change rate observed
\end{minipage} & \begin{minipage}[b]{\linewidth}\raggedright
corrected
\end{minipage} & \begin{minipage}[b]{\linewidth}\raggedright
cosmetic FP steps
\end{minipage} & \begin{minipage}[b]{\linewidth}\raggedright
steps
\end{minipage} \\
\midrule
DOM & 0.3846 & 0.4332 & 1618 & 33277 \\
phantom\_\allowbreak{}prompt & 0.4222 & 0.4681 & 1528 & 33283 \\
phantom\_\allowbreak{}som & 0.3683 & 0.4169 & 1688 & 34741 \\
phantom\_\allowbreak{}text & 0.3331 & 0.3770 & 1518 & 34610 \\
SoM & 0.4071 & 0.4675 & 1870 & 30921 \\
Vision & 0.5675 & 0.6134 & 1569 & 34151 \\
\bottomrule
\end{tabular}
\end{table}

% Table source-order #35: generated by scripts/analysis/export_ablation_tables.py
% Do not hand-edit the body; rerun the producer. Break-point fixes are the exception.
\begin{table}[tp]
\centering
\caption{Evaluator granularity. Evaluator granularity over the paper-grade set. Across 7,686 scored episodes in 36 conditions, the evaluator emits 2 distinct values (0.0, 1.0) (645 of them positive) with 0 missing and 0 non-numeric. There is no graded quality target to regress on. That is a property of the benchmark's design rather than of this pipeline, and it is a precondition of every routing negative in this document: a router's training signal can only be as fine-grained as the evaluator, so a partially-completed task is indistinguishable from one the agent never started. Source: \texttt{evaluator\_score\_granularity.json}.}
\label{tab:t35}
\footnotesize
\setlength{\tabcolsep}{3pt}
\begin{tabular}{@{}ll@{}}
\toprule
quantity & value \\
\midrule
conditions resolved & 36 \\
conditions unresolved & 0 \\
episodes scored & 7686 \\
episodes with no score & 0 \\
episodes with a non-numeric score & 0 \\
distinct score values & {[}0.0, 1.0{]} \\
count per value & \{``0.0'': 7041, ``1.0'': 645\} \\
\bottomrule
\end{tabular}
\end{table}


% Table source-order #28: Leaked-success sensitivity. Sensitivity to environmentally-credited su
\begin{table}[tp]
\centering
\caption{Leaked-success sensitivity. Sensitivity to environmentally-credited successes. \texttt{require\_reset} is a no-op on reddit, so subscriptions accumulate across a run's episodes and a later task can be scored on state an earlier one created. 6 such successes are set to 0 here: the denominator is unchanged, because an attempted-and-unaccomplished task is a 0, not a missing row. 4 of the leaks are on DOM, so removing them helps the fused arm: the direction that disfavours this project's own caution. \textbf{Caution:} The WA cells are unaudited for the same defect. Source: \texttt{leakage\_sensitivity.json}.}
\label{tab:t28}
\footnotesize
\setlength{\tabcolsep}{3pt}
\begin{tabular}{@{}
  >{\raggedright\arraybackslash}p{(\linewidth - 12\tabcolsep) * \real{0.1429}}
  >{\raggedright\arraybackslash}p{(\linewidth - 12\tabcolsep) * \real{0.1429}}
  >{\raggedright\arraybackslash}p{(\linewidth - 12\tabcolsep) * \real{0.1429}}
  >{\raggedright\arraybackslash}p{(\linewidth - 12\tabcolsep) * \real{0.1429}}
  >{\raggedright\arraybackslash}p{(\linewidth - 12\tabcolsep) * \real{0.1429}}
  >{\raggedright\arraybackslash}p{(\linewidth - 12\tabcolsep) * \real{0.1429}}
  >{\raggedright\arraybackslash}p{(\linewidth - 12\tabcolsep) * \real{0.1429}}@{}}
\toprule
\begin{minipage}[b]{\linewidth}\raggedright
cell
\end{minipage} & \begin{minipage}[b]{\linewidth}\raggedright
contrast
\end{minipage} & \begin{minipage}[b]{\linewidth}\raggedright
before
\end{minipage} & \begin{minipage}[b]{\linewidth}\raggedright
95\% CI
\end{minipage} & \begin{minipage}[b]{\linewidth}\raggedright
after
\end{minipage} & \begin{minipage}[b]{\linewidth}\raggedright
95\% CI
\end{minipage} & \begin{minipage}[b]{\linewidth}\raggedright
verdict
\end{minipage} \\
\midrule
red·B0 & SoM − Vision & +7.39 & {[}+2.46, +12.32{]} & +7.88 & {[}+2.96, +12.81{]} & unchanged \\
red·B0 & SoM − DOM & +0.49 & {[}-3.94, +4.93{]} & +0.99 & {[}-3.45, +5.42{]} & unchanged \\
red·B1 & SoM − Vision & +4.93 & {[}+1.48, +8.87{]} & +4.43 & {[}+0.99, +7.88{]} & unchanged \\
red·B1 & SoM − DOM & +1.48 & {[}-1.48, +4.43{]} & +0.99 & {[}-1.48, +3.94{]} & unchanged \\
red·B2 & SoM − Vision & -0.99 & {[}-3.45, +1.48{]} & -0.99 & {[}-3.45, +1.48{]} & unchanged \\
red·B2 & SoM − DOM & -2.96 & {[}-5.91, -0.49{]} & -1.48 & {[}-3.45, +0.49{]} & \textbf{flips} \\
\bottomrule
\end{tabular}
\end{table}

% Table source-order #40: Earned versus leaked successes. Which successes were earned. \texttt{\
\begin{table}[tp]
\centering
\caption{Earned versus leaked successes. Which successes were earned. \texttt{\#sidebar\ \textgreater{}\ section\ \textgreater{}\ ul} is read by 9 VWA reddit tasks; \texttt{require\_reset} is a no-op on reddit so subscriptions accumulate. LEAKED = scored success by an episode that never visited the required forum. On VWA: 6 leaked, 68 earned; Table~\ref{tab:t28} recomputes every contrast with the leaked ones zeroed. WebArena audited 2026-08-03 (first time): 50 scored episodes over its 5 sidebar tasks, 0 leaked. \textbf{Caution:} That zero is a \emph{lower bound}, not a clearance: the test asks whether the episode reached the forum, and an episode can arrive at a forum an earlier one subscribed to, read \texttt{Unsubscribe}, and finish without acting. One such case is hand-confirmed (\texttt{B1}/DOM task 597) and scores \texttt{earned} here; a text heuristic for the pattern was tried and rejected because model self-report cannot separate deliberating from acting. Source: \texttt{reddit\_sidebar\_leakage\_audit\_with\_wa.json}.}
\label{tab:t40}
\footnotesize
\setlength{\tabcolsep}{3pt}
\begin{tabular}{@{}lllll@{}}
\toprule
benchmark & cell · mode & scored successes & of which LEAKED & share \\
\midrule
VWA & B0 · DOM & 5 & 1 & 20.0\% \\
VWA & B0 · SoM-image & 2 & 0 & 0.0\% \\
VWA & B0 · DOM+\allowbreak{}sprompt & 3 & 0 & 0.0\% \\
VWA & B0 · DOM+\allowbreak{}stext & 3 & 0 & 0.0\% \\
VWA & B0 · SoM & 3 & 0 & 0.0\% \\
VWA & B0 · Vision & 2 & 1 & 50.0\% \\
VWA & B1 · DOM & 2 & 0 & 0.0\% \\
VWA & B1 · SoM-image & 4 & 0 & 0.0\% \\
VWA & B1 · DOM+\allowbreak{}sprompt & 2 & 0 & 0.0\% \\
VWA & B1 · DOM+\allowbreak{}stext & 4 & 0 & 0.0\% \\
VWA & B1 · SoM & 3 & 1 & 33.3\% \\
VWA & B2 · DOM & 3 & 3 & 100.0\% \\
VWA & B2 · SoM & 1 & 0 & 0.0\% \\
WA & B0 · DOM & 3 & 0 & 0.0\% \\
WA & B0 · SoM-image & 5 & 0 & 0.0\% \\
WA & B0 · DOM+\allowbreak{}stext & 5 & 0 & 0.0\% \\
WA & B0 · SoM & 3 & 0 & 0.0\% \\
WA & B0 · Vision & 3 & 0 & 0.0\% \\
WA & B1 · DOM & 4 & 0 & 0.0\% \\
WA & B1 · SoM-image & 5 & 0 & 0.0\% \\
WA & B1 · DOM+\allowbreak{}stext & 5 & 0 & 0.0\% \\
WA & B1 · SoM & 3 & 0 & 0.0\% \\
WA & B1 · Vision & 1 & 0 & 0.0\% \\
\bottomrule
\end{tabular}
\end{table}


% Paper B appendix, ported 2026-08-05 (merged-paper debt). Booktabs/width
% repairs applied; content otherwise verbatim from paperB_appendix.tex.
\section{Label definitions, nesting, and identifiability}\label{supporting-tables}

Every number in these tables is quoted in the body section that references it. They are placed here because the body has an eight-page limit and these tables document the results rather than making the argument.

\subsection{Observation modes}\label{observation-modes}

The grid \S\ref{sec:setup} refers to. The last column is the modality split Appendix~\ref{screenshot-tier} uses. It is not a cost split, for the reason \S\ref{sec:setup} gives: the image-bearing tier holds both the dearest mode and the cheapest.

\begin{table}[t]
\centering
\caption{The six observation modes (\S\ref{sec:setup}). The three P-modes keep either the mark legend, the SoM prompt, or both, while removing the per-step screenshot.}
\small
\setlength{\tabcolsep}{4pt}
\begin{tabular}{@{}llll@{}}
\toprule
mode & text payload & prompt family & annotated screenshot \\
\midrule
DOM & accessibility tree & DOM & no \\
SoM & mark legend & SoM & \textbf{yes} \\
Vision & none & vision & \textbf{yes} (unannotated) \\
P-text & mark legend & DOM & no \\
P-prompt & accessibility tree & SoM & no \\
P-SoM & mark legend & SoM & no \\
\bottomrule
\end{tabular}
\end{table}

\subsection{Where the which-mode label disagrees with measured cost}\label{where-the-which-mode-label-disagrees-with-measured-cost}

Per-cell detail for the 12.5--54.6\% range quoted in \S\ref{sec:lowerbound}, and for the claim that the exact-tie case never occurs.

\begin{table}[t]
\centering
\caption{Per-cell detail behind \S\ref{sec:lowerbound}. The last column counts labelled tasks on which the fixed priority list selected a mode strictly more expensive than another mode that also succeeded on that task. The exact-tie case, where the list order is the only tiebreaker, occurs in zero rows in every cell.}
\small
\setlength{\tabcolsep}{4pt}
\begin{tabular}{@{}llll@{}}
\toprule
cell & labels & multi-success & list picked a strictly pricier mode \\
\midrule
classifieds · B0 & 97 & 68 (70.1\%) & \textbf{53 (54.6\%)} \\
reddit · B0 & 53 & 36 (67.9\%) & \textbf{23 (43.4\%)} \\
classifieds · B1 & 55 & 29 (52.7\%) & \textbf{26 (47.3\%)} \\
reddit · B1 & 24 & 17 (70.8\%) & \textbf{9 (37.5\%)} \\
classifieds · B2 & 16 & 4 (25.0\%) & \textbf{2 (12.5\%)} \\
reddit · B2 & 15 & 3 (20.0\%) & \textbf{2 (13.3\%)} \\
\bottomrule
\end{tabular}
\end{table}

\subsection{Effect of nesting the operating-point selection}\label{effect-of-nesting-the-operating-point-selection}

Per-cell detail for the −0.99pp to +1.34pp range quoted in \S\ref{sec:lowerbound}, and for the observation that nesting moves the result in both directions.

\begin{table}[t]
\centering
\caption{Per-cell detail behind \S\ref{sec:lowerbound}. The naive column selects the threshold and the strongest and cheapest modes from whole-cell outcomes; the nested column re-derives all three inside each outer fold. Nesting moves the result in both directions.}
\small
\setlength{\tabcolsep}{4pt}
\begin{tabular}{@{}llll@{}}
\toprule
cell & naive nesting SR & fully nested SR & Δ \\
\midrule
classifieds · B0 & 25.45\% & 26.79\% & \textbf{+1.34pp} \\
classifieds · B1 & 13.84\% & 14.29\% & +0.45pp \\
classifieds · B2 & 1.34\% & 1.34\% & ±0.00pp \\
reddit · B0 & 13.79\% & 12.81\% & \textbf{−0.99pp} \\
reddit · B1 & 6.40\% & 5.91\% & −0.49pp \\
reddit · B2 & 3.94\% & 3.94\% & ±0.00pp \\
\bottomrule
\end{tabular}
\end{table}

\subsection{Identifiability of pooled which-mode labels}\label{identifiability-of-pooled-which-mode-labels}

Per-site detail for the conflict rates and modal-agreement figures quoted in Appendix~\ref{pooling-across-backbones}.

\begin{table}[t]
\centering
\caption{Per-site detail behind Appendix~\ref{pooling-across-backbones}. A conflict is one task on which two cells recorded different oracle modes. In-sample modal agreement is the accuracy of emitting the modal label per task, over all pooled labelled rows (168 on classifieds, 92 on reddit), a row being one (task, backbone) pair on which that backbone solved the task.}
\small
\setlength{\tabcolsep}{4pt}
\begin{tabular}{@{}
  >{\raggedright\arraybackslash}p{(\linewidth - 12\tabcolsep) * \real{0.1429}}
  >{\raggedright\arraybackslash}p{(\linewidth - 12\tabcolsep) * \real{0.1429}}
  >{\raggedright\arraybackslash}p{(\linewidth - 12\tabcolsep) * \real{0.1429}}
  >{\raggedright\arraybackslash}p{(\linewidth - 12\tabcolsep) * \real{0.1429}}
  >{\raggedright\arraybackslash}p{(\linewidth - 12\tabcolsep) * \real{0.1429}}
  >{\raggedright\arraybackslash}p{(\linewidth - 12\tabcolsep) * \real{0.1429}}
  >{\raggedright\arraybackslash}p{(\linewidth - 12\tabcolsep) * \real{0.1429}}@{}}
\toprule
\begin{minipage}[b]{\linewidth}\raggedright
site
\end{minipage} & \begin{minipage}[b]{\linewidth}\raggedright
tasks labelled in ≥2 cells
\end{minipage} & \begin{minipage}[b]{\linewidth}\raggedright
conflicting
\end{minipage} & \begin{minipage}[b]{\linewidth}\raggedright
conflict rate
\end{minipage} & \begin{minipage}[b]{\linewidth}\raggedright
in-sample modal agreement
\end{minipage} & \begin{minipage}[b]{\linewidth}\raggedright
same, on shared tasks only
\end{minipage} & \begin{minipage}[b]{\linewidth}\raggedright
same, exact-vector grouping
\end{minipage} \\
\midrule
classifieds & 54 & 31 & \textbf{57.4\%} & \textbf{79.2\%} & 70.3\% (n=118) & 83.9\% \\
reddit & 25 & 14 & \textbf{56.0\%} & \textbf{83.7\%} & 74.1\% (n=58) & 89.1\% \\
\bottomrule
\end{tabular}
\end{table}

\emph{We call it agreement and not a Bayes ceiling because it is a resubstitution estimate: it scores the same rows it took the modal label from, so a task labelled by only one backbone is correct by construction. That describes 50 of 168 classifieds rows (29.8\%) and 34 of 92 reddit rows (37.0\%); restricted to tasks two or more backbones label, agreement falls to 70.3\% and 74.1\%. An out-of-sample bound would need leave-one-backbone-out prediction or a shrinkage estimator, and would be lower still. Every number here is therefore an optimistic bound on what a pooled classifier could reach, which is the direction the argument needs.}

\emph{The last column groups by the exact feature vector instead of by task. Rows of one task are not always identical, because three of the five numeric features are read from that backbone's own step-0 observation, so they differ somewhere on 31.5\% of shared classifieds tasks and 80.0\% of shared reddit tasks. We report but do not use that grouping: it leaves 74 of 117 classifieds groups and 69 of 78 reddit groups with a single member, covering 44\% and 75\% of rows, so it is even more inflated than the headline, and a router serving one backbone could not recover backbone identity from that jitter anyway.}

\subsection{The screenshot-modality tier}\label{the-screenshot-modality-tier}

Per-site detail for Appendix~\ref{screenshot-tier}, on the same pooled labelled rows and the same grouping as Appendix~\ref{identifiability-of-pooled-which-mode-labels}. The agreement column is the exception: it is defined only on tasks labelled by two or more backbones, since agreement needs two labels to compare.

\begin{table}[t]
\centering
\caption{Per-site detail behind Appendix~\ref{screenshot-tier}, over the same solve events. Columns two and three carry the same resubstitution caveat as Appendix~\ref{identifiability-of-pooled-which-mode-labels}, and column three additionally rises for an arithmetic reason: merging six classes into two can only increase a modal share. The claim that backbones agree about the screenshot therefore rests on the last two columns, which measure agreement between two backbones' labels directly rather than against a modal label. Under Appendix~\ref{identifiability-of-pooled-which-mode-labels}'s exact-vector grouping the tier figures are 92.3\% and 97.8\%. No classifier is fitted to this target anywhere in the paper.}
\small
\setlength{\tabcolsep}{4pt}
\begin{tabular}{@{}
  >{\raggedright\arraybackslash}p{(\linewidth - 8\tabcolsep) * \real{0.2000}}
  >{\raggedright\arraybackslash}p{(\linewidth - 8\tabcolsep) * \real{0.2000}}
  >{\raggedright\arraybackslash}p{(\linewidth - 8\tabcolsep) * \real{0.2000}}
  >{\raggedright\arraybackslash}p{(\linewidth - 8\tabcolsep) * \real{0.2000}}
  >{\raggedright\arraybackslash}p{(\linewidth - 8\tabcolsep) * \real{0.2000}}@{}}
\toprule
\begin{minipage}[b]{\linewidth}\raggedright
site
\end{minipage} & \begin{minipage}[b]{\linewidth}\raggedright
which-mode modal agreement
\end{minipage} & \begin{minipage}[b]{\linewidth}\raggedright
tier modal agreement
\end{minipage} & \begin{minipage}[b]{\linewidth}\raggedright
tier agreement across backbones
\end{minipage} & \begin{minipage}[b]{\linewidth}\raggedright
six-way agreement across backbones
\end{minipage} \\
\midrule
classifieds & 79.2\% & \textbf{89.9\%} & \textbf{68.5\%} & 42.6\% \\
reddit & 83.7\% & \textbf{96.7\%} & \textbf{88.0\%} & 44.0\% \\
\bottomrule
\end{tabular}
\end{table}

\section{Derivations for the four relabelling routes}\label{derivations-for-the-four-relabelling-routes}

\S\ref{sec:lowerbound} and \S\ref{sec:gap} state the outcome of each route. This appendix gives the derivation.

\subsection{Continuous labels}\label{continuous-labels}

The cleanest fix would be to regress on a graded quality signal rather than classify a discrete winner: partial credit turns every episode into a training example regardless of whether it succeeded. VisualWebArena does not provide one. Across the 7,686 scored episodes of our 36 landed conditions the evaluator emits exactly two values, 0.0 and 1.0, in a 7,041 / 645 split, as do the 36 protocol-excluded episodes we do not score. This is a property of the benchmark's evaluation design rather than of our pipeline, and it forecloses the route entirely.

\subsection{Coarser classes}\label{coarser-classes}

\S\ref{sec:lowerbound}'s obstruction is that four of six cells have fewer than ten labelled rows in more than one class. Merging classes does not add rows. The binary collapse of the screenshot tier (Appendix~\ref{screenshot-tier}) does help, but for a different reason, having to do with agreement across backbones rather than with class count.

\subsection{Pooling across backbones}\label{pooling-across-backbones}

Six cells at 15--97 labels become 260 pooled examples and every class clears the minimum-count filter, so supply is solved. Identifiability is not. The features carry no model identity, so two backbones facing the same task on the same site produce near-identical feature vectors, and where they are identical and the oracle labels differ, a classifier is being asked to emit two different answers for one input. The figure reported in Appendix~\ref{identifiability-of-pooled-which-mode-labels} is the accuracy of emitting the modal label per group, scored on the rows the label came from. It is an in-sample bound on what a pooled classifier could reach, not a Bayes ceiling, and Appendix~\ref{identifiability-of-pooled-which-mode-labels} gives the resubstitution caveat.

They are only near-identical, not identical, because \texttt{dom\_\allowbreak{}complexity}, \texttt{text\_\allowbreak{}length} and \texttt{tokens\_\allowbreak{}input\_\allowbreak{}text} are read from the backbone's own step-0 observation rather than from the task config. On 31.5\% of shared classifieds tasks and 80.0\% of shared reddit tasks the rows therefore differ somewhere. Grouping by the exact vector rather than by task raises the figure (Appendix~\ref{identifiability-of-pooled-which-mode-labels}, last column), and we report but do not adopt that number: it leaves most groups with one member, and a group of one is scored perfectly whatever the labels do, so it inflates with feature sparsity rather than tracking identifiability. Every version of the number is an optimistic bound and all of them are far below what a deployable which-mode router needs, which is the only thing the argument turns on.

\subsection{Screenshot tier}\label{screenshot-tier}

The tier label is derived from the same solve events as the which-mode label, by mapping each oracle mode to image-bearing (SoM, Vision) or text-only (DOM, P-text, P-prompt, P-SoM). No new episodes are involved, which is why the figure rises without a single new solve event. Part of that rise is arithmetic (two classes admit a larger modal share than six), so the agreement columns rather than the modal-agreement columns carry the claim. The tier is defined only on tasks that some mode solved, so its denominator is the solved set and not the full task universe. Appendix~\ref{the-screenshot-modality-tier}'s modal-agreement columns run over the whole pooled labelled set, as in Appendix~\ref{identifiability-of-pooled-which-mode-labels}; the two agreement columns are restricted to tasks labelled by two or more backbones, that being the only set on which cross-backbone agreement is defined.

\section{The reddit · B2 saving in detail}\label{the-reddit-b2-saving-in-detail}

\S\ref{sec:lowerbound} reports that the one cell whose cost saving survives Holm correction has an AUROC below chance. The mechanism is tail enrichment rather than a globally ordered score.

Reddit · B2 sends 192 of 203 tasks (95\%) to the cheap mode with no accuracy loss, which is unsurprising in a cell where only 7.4\% of tasks are solvable at all: almost nothing in that 95\% was going to succeed under any mode. The policy therefore differs from the free always-cheapest policy by five percent of the allocation. The 11 tasks it holds back for the strong mode carry four successes that the fixed policy does not collect, against four collected by the fixed policy overall. The permutation null detects that enrichment. A globally ordered score is not required to produce it, which is why the cell's AUROC of 0.483, below both chance and its own best single covariate at 0.711, is consistent with a real saving.

Two properties of that test are worth recording. It runs 10,000 draws and reports the plus-one Monte Carlo estimator (k+1)/(B+1), whose floor is therefore 1.0 × 10⁻⁴, two orders below the tightest Holm threshold of 8.3 × 10⁻³. That matters because at the 200 draws we first used, the floor was 1/201 = 4.98 × 10⁻³, this cell reported exactly it, and whether it could clear the threshold at all was a function of the draw count rather than of the data; at 10,000 draws four of the draws match or beat the observed saving, so p = 5.0 × 10⁻⁴ is measured. Second, the quantity tested is the saving at an operating point selected against whole-cell outcomes, which is not the nested policy of Table 6. Null and observation select that point the same way, so the null absorbs the selection optimism and the comparison is fair, but the point is not one a deployment could occupy, and \S\ref{sec:lowerbound}'s conclusion rests on the nested numbers rather than on this test.

\subsection{Supply and trainability under both label definitions}\label{supply-and-trainability-under-both-label-definitions}

\S\ref{sec:lowerbound} keeps the prior-order label and reports the measured-cost rule as a sensitivity. Supply is identical under both by construction, since each labels exactly the tasks some mode solved, so only the class distribution and through it trainability can move. The relabelled column equals the ``order picked a strictly pricier mode'' column of A.2 exactly, which is the consistency check: a label moves if and only if the order's pick was not the measured cheapest.

\begin{table}[t]
\centering
\caption{Trainability under the two label definitions. ``Surviving'' counts classes clearing ten training rows in a five-fold split; \textbf{no} marks a cell with fewer than two, where no classifier exists. Four of six cells are untrainable under the reported label and five of six under the measured-cost alternative, so the supply argument does not depend on the choice.}
\small
\setlength{\tabcolsep}{4pt}
\begin{tabular}{@{}
  >{\raggedright\arraybackslash}p{(\linewidth - 8\tabcolsep) * \real{0.2000}}
  >{\raggedright\arraybackslash}p{(\linewidth - 8\tabcolsep) * \real{0.2000}}
  >{\raggedright\arraybackslash}p{(\linewidth - 8\tabcolsep) * \real{0.2000}}
  >{\raggedright\arraybackslash}p{(\linewidth - 8\tabcolsep) * \real{0.2000}}
  >{\raggedright\arraybackslash}p{(\linewidth - 8\tabcolsep) * \real{0.2000}}@{}}
\toprule
\begin{minipage}[b]{\linewidth}\raggedright
cell
\end{minipage} & \begin{minipage}[b]{\linewidth}\raggedright
labels
\end{minipage} & \begin{minipage}[b]{\linewidth}\raggedright
relabelled
\end{minipage} & \begin{minipage}[b]{\linewidth}\raggedright
prior order: surviving classes
\end{minipage} & \begin{minipage}[b]{\linewidth}\raggedright
measured cost: surviving classes
\end{minipage} \\
\midrule
classifieds · B0 & 97 & 53 & 3 (DOM, P-prompt, SoM) & 2 (SoM, Vision) \\
reddit · B0 & 53 & 23 & 1 (DOM) (\textbf{no}) & 1 (P-text) (\textbf{no}) \\
classifieds · B1 & 55 & 26 & 2 (DOM, SoM) & 1 (Vision) (\textbf{no}) \\
reddit · B1 & 24 & 9 & 0 (\textbf{no}) & 0 (\textbf{no}) \\
classifieds · B2 & 16 & 2 & 0 (\textbf{no}) & 0 (\textbf{no}) \\
reddit · B2 & 15 & 2 & 0 (\textbf{no}) & 0 (\textbf{no}) \\
\bottomrule
\end{tabular}
\end{table}

The single cell that changes, classifieds · B1, loses a class rather than gaining one: the prior-order label keeps DOM and SoM above the threshold, the measured-cost label concentrates enough of those rows onto Vision that only Vision survives. The alternative definition therefore strengthens the negative result, which is a reason to report it and not a reason to adopt it.

\subsection{The best-success mode is not stable across folds}\label{the-best-success-mode-is-not-stable-across-folds}

The nested design of \S\ref{sec:lowerbound} re-selects the best-success mode inside every outer fold, which exposes something the whole-cell version conceals. In reddit · B0 the five outer folds select DOM, DOM, SoM, SoM, DOM. A pipeline that picks one best mode from all realised outcomes is therefore not merely optimistic about its threshold; it reports a mode choice that its own resampling does not reproduce.

